% concatenated from homgo_en.tex
\documentclass[preprint]{elsarticle}
%%% begin file('ru//in/head.tex') %%%
\usepackage{cmap} % for russuan cope paste
           \usepackage[T2A]{fontenc}   % Or any other encoding.
           \usepackage[utf8]{inputenc}
           \usepackage[russian,english]{babel}
\usepackage{amsmath}
\usepackage{amsfonts}
\usepackage{amssymb}
\usepackage{amsthm}
\usepackage[pdftex,unicode]{hyperref}
\usepackage{indentfirst} % включить отступ у первого абзаца
%\usepackage{amscd}
%\usepackage[all,cmtip]{xy}
%\usepackage{mathtools}
%\usepackage{comment}
\usepackage{tikz-cd}
% !!!!!
%\usetikzlibrary{babel}

%\usepackage[inline]{showlabels}

\def\Sl{\mathbb S} \def\J{\mathbb J} \def\R{\mathbb R}
\def\T{\mathbb T}
\def\N{\mathbb N}    
\def\C{\mathbb C}  


\def\uu{\Upsilon}
\def\vu{\Psi}
\def\UU{\mathfrak{U}}

\def\om{\omega}
\def\Om{\Omega}
\def\aut#1{\operatorname{Aut}(#1)}
\def\supp{\operatorname{supp}}
\def\autp#1#2{\operatorname{Aut_{#1}}(#2)}
%\def\be{\operatorname{\beta}}
\def\be{\beta}
\def\dim{\operatorname{dim}}
\def\al{\alpha}
\def\ga{\gamma}
\def\la{\lambda}
\def\ph{\varphi}
\def\de{\delta}
\def\ka{\kappa}

\def\lll{\Lambda}

\def\ggg{\Gamma}
\def\gbm{\ggg^{\scriptscriptstyle{BM}}}
\def\god{\ggg^{\scriptscriptstyle{OD}}}

\def\ode{{\wt{OD}}}
\def\bme{{\wt{BM}}}

\def\nfs/{NFS}
\def\cdp/{CDP}
\def\cdpz/{CDP${}_0$}
\def\cnL#1{$(L_{#1})$}
\def\cnLl#1{$(L'_{#1})$}

\def\lleq{\preccurlyeq}
\def\ggeq{\succcurlyeq}

% mathabx
% mathb

\def\cQ{\mathcal{Q}}


\def\cA{\mathcal A}
\def\cB{\mathcal{B}}
\def\cE{\mathcal{E}}
\def\cW{\mathcal{W}}
\def\cN{\mathcal{N}}
\def\cL{\mathcal{L}}
\def\cH{\mathcal{H}}
\def\cp{\mathcal{p}}
\def\cp{\mathcal{S}}
\def\s{{\mathfrak{Q}}}
\def\cG{\mathcal{G}}
\def\gG{{\mathfrak{G}}}

\def\gr{{\mathfrak{R}}}

\def\ll{{\mathfrak{L}}}
\def\uu{\Upsilon}
\def\mm{{\mathfrak{M}}}
\def\ww{{\mathfrak{W}}}
%\def\vv{{\mathfrak{V}}}
\def\ppi{\varpi}






\def\cl#1{\overline{#1}}
\def\clx#1#2{\overline{#2}^{\,#1}}
\def\clv#1{\clx{v}{#1}}
%\def\int#1{\operatorname{int} (#1)}
%\def\intx#1#2{\operatorname{int}_{#1} (#2)}
\def\Int{\operatorname{Int}}

\def\id{\operatorname{id}}
\def\sp{\operatorname{sp}}
\def\bt{\operatorname{\beta}}

\def\ltu{\operatorname{\cl{lt}}}
%\DeclareMathOperator*{\ltu}{\cl{lt}}
\def\oms{\om^*}

\def\limp#1{{\textstyle \lim_{#1}}}

\def\es{\varnothing}
\def\Tau{{\mathcal T}}


\def\sset#1{\{#1\}}
\def\fset#1{\{#1\}}
\def\set#1{\bbset#1\eeset}
\def\bbset#1:#2\eeset{\{#1\,:\,#2\}}
\def\sett#1{\bbsett#1\eesett}
\def\bbsett#1:#2\eesett{\{#1\,:\,\text{#2}\}}

\def\iset#1{\ibbset#1\ieeset}
\def\ibbset#1:#2\ieeset{(#1)_{#2}}


\def\go{{\mathfrak{O}}} 
\def\gs{{\mathfrak{S}}}
\def\gm{{\mathfrak{M}}}
\def\cm{{\mathcal M}}
\def\cM{{\mathcal M}}



\def\tp{\Tau}
\def\tps{\tp^*}
\def\tpss#1{\tps(#1)}
\def\ppi#1{\Pi(#1)}
\def\cP{{\mathcal P}}
\def\cB{{\mathcal B}}
\def\gP{{\mathfrak{P}}}

\def\pbase#1{\gP\left[#1\right]}


\def\gF{{\mathfrak{F}}}
\def\gB{{\mathfrak{B}}}
\def\cB{{\mathcal B}}
\def\cF{{\mathcal F}}
\def\cS{{\mathcal S}}
\def\cU{{\mathcal U}}
\def\cO{{\mathcal O}}
\def\cR{{\mathcal R}}
\def\cD{{\mathcal D}}



\def\Exp#1{\operatorname{Exp}(#1)}
\def\Expne#1{\operatorname{{Exp}}_*(#1)}
\def\ult#1{\operatorname{Ult}(#1)}
\def\ultf#1{\operatorname{Ult}_f(#1)}
\def\ultrf#1{\ultf{#1}}
%\def\filtr#1{\gF(#1)}
\def\filtr#1{\operatorname{Filtr}(#1)}
\def\centr#1{\operatorname{Centr}(#1)}


%\def\g{{\mathsf{G}}}

\def\gi{{\mathfrak{i}}}
\def\gm{{\mathfrak{m}}}



\def\til{\tilde}
\def\tD{{\til{\D}}}
% widetilde
\def\tV{\widetilde{V}}
\def\tM{\widetilde{M}}
\def\wt#1{\widetilde{#1}}
\def\tx{{\til{x}}}
\def\ts{{\til{s}}}
\def\tS{{\wt S}}
%\def\S{\mathbb S}
\def\gp{{\mathfrak{p}}}
\def\eqdef{\coloneqq}
\def\Ss{{\mathcal S}}
\def\ci{{\mathcal i}}
\def\cI{{\mathcal I}}
\def\cX{{\mathcal X}}
\def\cV{{\mathcal V}}
\def\cN{{\mathcal N}}



\def\ct{\tau}
\def\cT{{\mathcal T}}
\def\cs{{\mathcal s}}
\def\Mu{\mathrm{M}}


\def\E{{\sset\es}}

\def\diag{\mathop{\bigtriangleup}}

\def\nov#1{\operatorname{Nov}(#1)}
\def\dev#1{\operatorname{dev}(#1)}
\def\St{\operatorname{St}}
\def\st{\operatorname{st}}

%\def\fun{\times}
%\def\fun{\cup}
\def\fun{{}^\frown}

\def\AND{\bigwedge}
\def\OR{\bigvee}

\newcommand\restrA[2]{{% we make the whole thing an ordinary symbol
  \left.\kern-\nulldelimiterspace % automatically resize the bar with \right
  #1 % the function
  \vphantom{\big|} % pretend it's a little taller at normal size
  \right|_{#2} % this is the delimiter
  }}

\newcommand\restrB[2]{\ensuremath{\left.#1\right|_{#2}}}

\def\restr#1#2{\restrB{#1}{#2}}

\def\te{\theta}


\def\D{\Delta}
\def\term#1{{\it #1}}


% 15.02.2022 13:42:59 
\def\wh#1{\widehat{#1}}
%\def\tg#1#2{\wh{#1#2}}
\def\et(#1){ (#1)}

%\newtheorem{assertion}{Statement}
\newtheorem{statement}{Statement}[section]
\newtheorem{claim}{Fact}
\newtheorem{proposition}{Proposition}[section]
\newtheorem{theorem}{Theorem}[section]
\newtheorem*{theorem*}{Theorem}
\newtheorem{lemma}{Lemma}
\newtheorem*{lemma*}{Lemma}
\newtheorem{cor}{Corollary}[section]

\theoremstyle{definition}
\newtheorem{example}{Example}
%[section]

\newtheorem{definition}{Definition}[section]
\newtheorem{question}{Question}
%[section]
\newtheorem{task}{Task}[section]
\newtheorem{problem}{Problem}[section]
\newtheorem{hypothesis}{Hypothesis}[section]
\theoremstyle{remark}
\newtheorem{note}{Remark}[section]
\newtheorem*{note*}{Remark}

%%%%%%%%%%%%%%%%%%%%%%%%%%%%%%

\def\ep{\varepsilon}


%%%%%%%%%%%%%%%%%%%%%%%%



\def\gd/{$G_\delta$}
\def\F{\Phi}


\def\rarr{\Rightarrow}
\def\larr{\Leftarrow}
\def\lrarr{\Leftrightarrow}


%\def\lrar{\leftrightarrow}
\def\La{\Lambda}

\def\gm{{\mathfrak{m}}}
\def\bx{\bar x}
\def\wh#1{\widehat{#1}}
\def\hF{\wh\F}
\def\ii{{\mathfrak{I}}}

\def\DD{\operatornamewithlimits{%
  \mathchoice{\vcenter{\hbox{\huge \Delta}}}
             {\vcenter{\hbox{\Large \Delta}}}
             {\mathrm{\Delta}}
             {\mathrm{\Delta}}}}
             
\def\DD{\operatornamewithlimits{\mbox{\hbox{\huge \Delta}}}}

\def\DD{\operatornamewithlimits{\Delta}}

\def\DD{\operatornamewithlimits{%
  \mathchoice{\vcenter{\hbox{\huge $\Delta$}}}
             {\vcenter{\hbox{\Large $\Delta$}}}
             {\mathrm{\Delta}}
             {\mathrm{\Delta}}}}
             

\def\DD{\operatornamewithlimits{%
  \mathchoice{\vcenter{\hbox{\huge $\bigtriangleup$}}}
             {\vcenter{\hbox{\Large $\bigtriangleup$}}}
             {\mathrm{\bigtriangleup}}
             {\mathrm{\bigtriangleup}}}}


       
             
\def\diag{\DD}
\def\diagsmall{\bigtriangleup}

\def\mm{{\mathcal{M}}}
\def\si{\sigma}
\def\tht{\theta}
\def\pch{\pi\chi}

\def\vh#1{\hat{#1}}

\def\gF{{\mathfrak{F}}}


\def\vt#1{\tilde{#1}}

\def\rlim{\varprojlim}

\renewcommand{\labelenumi}{{\rm(\arabic{enumi})}}
\renewcommand{\labelenumii}{{\rm(\alph{enumii})}}

%\def\lex#1#2{\mathfrak{L}(#2,#1)}
%\def\blex#1#2#3#4{\mathfrak{B}(#1,#2;#4,#3)}

\def\Q{\mathbb Q}
\def\Z{\mathbb Z}
\def\cf{\operatorname{cf}}
\def\rev{\operatorname{rev}}


\def\lex#1#2{\mathfrak{L}(#2^{#1})}
\def\lex#1#2{#2^{#1}}



\def\blex#1#2#3#4{\mathfrak{B}(#1,#2;#4^{#3})}
\def\blex#1#2#3#4{\mathfrak{B}_{#4^{#3}}(#1,#2)}
\def\blex#1#2#3#4{\mathfrak{B}_{#4}(#1,#2)}

\def\dom{\operatorname{Dom}}

\def\homa#1#2{\mathfrak{D}_1(#1,#2)}

\def\homg#1,#2#3{\mathfrak{D}_{#1}(#2,#3)}
\def\homm#1{{\mathfrak{D}(#1)}}


\def\chars#1{\chi_{#1}}


\def\stab#1{\Lambda(#1)}
\def\staba#1#2{\Theta(#1,#2)}
\def\stabf#1{\theta(#1)}
\def\stabl#1{\zeta(#1)}
\def\stablp#1#2{\zeta_{#1}(#2)}


\def\stab#1{\Lambda_{#1}}
\def\staba#1#2{\Theta(#1,#2)}
\def\stabf#1{\theta_{#1}}
\def\stabl#1{\phi_{#1}}
\def\stablp#1#2{\phi_{#2,#1}}


\def\dd#1#2{\de(#1,#2)}

% Latex string concatenation notation
\def\cnc#1#2{#1\oplus#2}
\def\cnc#1#2{{#1}^\frown#2}
\def\symblecnc{\mathbin{{}^\frown}}
\def\cnc#1#2{{#1}\symblecnc{#2}}

\def\supp#1{#1^{-1}(1)}


\def\wob/{$WOB$}

\def\lL(#1,#2){(-\infty,#1)}
\def\lR(#1,#2){(#1,+\infty)}


%\def\B{\mathbb D}
\def\B{2}
\def\B{\mathsf{2}}
\def\B{\mathbf{2}}

%\def\cj/{{\rm (J)}}
\def\cjx/#1,#2,#3/{{\rm (J${}_{#2,#3}^{#1}$)}}
\def\cj/#1,#2/{{\rm (J${}_{#1,#2}$)}}

\def\cjx/#1,#2,#3/{{\rm (${\mathrm J}_{#2,#3}^{#1}$)}}
\def\cj/#1,#2/{\cjx/1,#1,#2/}

\def\cJ{{\mathcal J}}

\def\ilcomma/{\boldsymbol{,}}
\def\ilsymlp{\boldsymbol{(}}
\def\ilsymrp{\boldsymbol{)}}
\def\ilsymlc{\boldsymbol{[}}
\def\ilsymrc{\boldsymbol{]}}

\def\ilfont#1{\boldsymbol{#1}}

\def\ilfontX#1{\mathrm{\text{\texttt{\textbf{\large #1}}}}}
\def\ilfontX#1{{\text{\texttt{\textbf{\large \upshape #1}}}}}

\def\ilfont#1{\raisebox{-1.2pt}{\ilfontX{#1}}}
\def\ilfontC#1{\raisebox{-0.7pt}{\ilfontX{#1}}}

\def\ilcomma/{\ilfontC{,}}
\def\ilsymlp{\ilfont{(}}
\def\ilsymrp{\ilfont{)}}
\def\ilsymlc{\ilfont{[}\,}
\def\ilsymrc{\,\ilfont{]}}


\def\ilpp(#1,#2){\ilsymlp#1\ilcomma/ #2\ilsymrp}
\def\ilcp[#1,#2){\ilsymlc#1\ilcomma/ #2\ilsymrp}
\def\ilpc(#1,#2]{\ilsymlp#1\ilcomma/ #2\ilsymrc}
\def\ilcc[#1,#2]{\ilsymlc#1\ilcomma/ #2\ilsymrc}

\def\lL(#1,#2){\ilpp(-\infty, #1)}
\def\lR(#1,#2){\ilpp(#1, +\infty)}

\def\jumps#1{\mathfrak{J}(#1)}

\def\nom{{n\in\om}}
\def\sq#1#2{(#1)_{#2}}
\def\sqn#1{\sq{#1}\nom}
\def\sqnn#1{\sqn{#1_n}}

\def\lexab{\lex\la\B}

\def\sS{\mathbf{S}}
\def\Rs{\mathbf{R}}
\def\Hs{\mathbf{H}}

\def\wS{\wt\sS}
\def\wcS{\wt\cS}

\def\oip#1{\cH_{+}(#1)}
\def\oim#1{\cH_{-}(#1)}
\def\oipm#1{\cH_{\pm}(#1)}

\def\mns{0_{\om_1}}
\def\mxs{1_{\om_1}}

\def\nr{N_r}
\def\nl{N_l}
\def\na{N_*}

\def\er{\cB_r}
\def\el{\cB_l}
\def\ea{\cB_*}

\def\ex{\cE}
\def\wex{\wt\cE}


%%% end file('ru//in/head.tex') %%%

\begin{document}
%\maketitle

\begin{frontmatter}

\title{Homogeneous linearly ordered spaces}
%\author{Evgenii Reznichenko} 

\author{Anton Lipin}
\ead{tony.lipin@yandex.ru}
\address{N.N.\ Krasovskii Institute of Mathematics and Mechanics, Ural Branch of the Russian Academy of Sciences, Yekaterinburg, Russia}

\author{Evgenii\corref{cor} Reznichenko}
\ead{erezn@inbox.ru}
\address{Department of General Topology and Geometry, Mechanics and  Mathematics Faculty, M.~V.~Lomonosov Moscow State University, Leninskie Gory 1, Moscow, 199991 Russia}


\cortext[cor]{Corresponding author}

%{\it Ural Federal University, Yekaterinburg, Russia}

%\bigskip



\begin{abstract}
Every compact subset of a homogeneous generalized ordered (GO) space has character at most $\om_1$ and cardinality at most $2^{\om_1}$; if such a subset has uncountable character, then the character of the whole space equals $\om_1$ and its $\pi$-character is countable.
We construct a homogeneous $\si$-compact linearly ordered space (LOTS) $\Hs$ containing a compact subset $\sS$ of cardinality $2^{\om_1}$ whose character is $\om_1$ at every point and whose weight and Souslin number are both $2^{\om_1}$; thus both bounds obtained are sharp.
We prove that a semitopological group that is a GO space is hereditarily paracompact; if, in addition, it is not a $P$-space, then it is submetrizable, has countable character, and its compact subsets are metrizable.
Every linearly ordered semitopological group (and, more generally, every GO semitopological group) is either metrizable or is a $P$-space; the same holds for topological groups.
We also show that in an order-homogeneous LOTS every compact subset is first countable.
\end{abstract}
\begin{keyword}
linearly ordered spaces
\sep
generalized ordered spaces
\sep
homogeneous spaces
\sep
compact spaces
\end{keyword}
\end{frontmatter}


\section{Introduction} \label{sec:intro}
%%% begin file('ru//sec/intro.tex') %%%

A space $X$ is called \term{homogeneous} if for all $x,y\in X$ there exists a self-homeomorphism $f: X\to X$ such that $f(x)=y$.
The theory of homogeneous spaces is a large and active area of general topology \cite{Arhangelskii_vanMill}. Compact homogeneous spaces are of particular interest, but the structure of compact subsets of homogeneous spaces is also studied. It is proved in \cite{arh67,Arhangelskii2018,rezn2020} that compact subsets of an extremally disconnected space are finite. In \cite{rezn2020,Groznova2023ru} this result is extended to homogeneous subspaces of finite products of extremally disconnected spaces.
Under the continuum hypothesis ($\mathrm{CH}$), homogeneous subspaces of countable products of extremally disconnected spaces are metrizable \cite{rezn2020}.

In this paper we study linearly ordered topological spaces (LOTS), their subspaces (GO spaces), and topological properties of homogeneous LOTS and GO spaces.
A homogeneous GO space of pointwise countable type has countable character \cite[Corollary 29]{Arhangelskii2005}.
To establish this result, it suffices to show that compact subsets of such spaces have countable character.

The main technical tool of the paper is a characterization of points of a GO space via linearly ordered neighbourhood bases. We call a point $x$ of a space $X$ a \wob/-point if it has a neighbourhood base that is linearly ordered by inclusion; for GO spaces this condition turns out to be equivalent to the character and the $\pi$-character coinciding at that point (Proposition \ref{p:1+2}). Together with the spectrum $\sp_t(x,X)$ of the left and right tightnesses at a point, this characterization reduces the computation of cardinal invariants of homogeneous GO spaces to the analysis of the tightnesses $\cf_l$, $\cf_r$ (Section \ref{sec:cfgos}), which is what allows us to obtain the main results of the paper.

In a homogeneous GO space $X$, every compact subset has character at most $\om_1$, and if $X$ contains a compact subset of uncountable character, then the character of $X$ equals $\om_1$ and its $\pi$-character is countable (Theorem \ref{t:main:1}, Section \ref{sec:hcomp}). Both bounds in this theorem are sharp. There exists a homogeneous countably compact LOTS $X_\la$ (for a regular ordinal $\la$) with $\chi(X_\la)=\la$ and $\pi\chi(X_\la)=\om$ (Example \ref{e:3}); for $\la=\om_1$, the space $X_{\om_1}$ additionally embeds the transfinite line $\om_1+1$ (Example \ref{e:4}). The bound on the cardinality of a compact subset is also attained: there exists a homogeneous $\si$-compact LOTS $\Hs$ containing a compact subset $\sS$ whose character at every point equals $\om_1$, while its weight, Souslin number and cardinality all equal $2^{\om_1}$ (Example \ref{e:5}).

A separate line of results concerns the stronger condition of \emph{order-homogeneity} of a space (when homogeneity is realized by self-homeomorphisms that preserve or reverse the order). It is proved that in an order-homogeneous LOTS every compact subset is first countable (Theorem \ref{t:cssoh:1}); this is substantially stronger than the bound $\chi(K)\leq \om_1$ valid for merely homogeneous spaces.

Let $G$ be a semitopological group that is a GO space. Then $G$ is hereditarily paracompact (Theorem \ref{t:main:1+1}, which strengthens \cite[Corollary 2.7.]{BuzyakovaVural2014}), and compact subsets of $G$ are metrizable (Corollary \ref{с:main:1}).
If, in addition, $G$ is not a $P$-space, then $G$ is a first-countable submetrizable space (Theorem \ref{t:main:2}).
If, in addition, $G$ is a topological group, or $G$ is a LOTS semitopological group, then $G$ is metrizable (Theorem \ref{t:main:2+1} and Theorem \ref{t:main:3}).

This result cannot be generalized to monotonically normal groups --- there are many examples of separable, non-first-countable, non-metrizable monotonically normal topological groups \cite{GartsideReznichenko2000nmfs}. Monotonically normal paratopological groups are hereditarily paracompact \cite[Theorem 2.6]{BuzyakovaVural2014}, and hence compact subsets of such groups have countable tightness \cite[Theorem 3.5]{WilliamsZhou1998} (Theorem \ref{t:main:3+1}). Compact subsets of monotonically normal topological groups are metrizable \cite[Theorem 2a]{Gartside1998}.
%(Theorem \ref{t:main:3+2}).
LOTS topological groups can have arbitrarily large character (Proposition \ref{p:h:tg1}).

%, and if $G$ is a $P$-space, then $G$ is a non-archimedean space

Papers \cite{BennettBurkeLutzer2013} considers GO and LOTS spaces equipped with algebraic structures weaker than a group structure.

Section \ref{sec:qe} discusses a number of open questions arising from our results: whether the LOTS $X_{\om_1}$ and $\Hs$ admit the structure of a right-topological group (by analogy with the well-known example of the compact LOTS `double arrow' \cite{Banakh2008}); whether the tightness of compact subsets of power homogeneous (monotonically normal) GO spaces is bounded by $\om_1$; and whether the theorems on hereditary paracompactness, the stratifiable property, and metrizability of compacta known for monotonically normal topological and paratopological groups extend to semitopological and quasitopological groups (in connection with the questions of Collins \cite[Problem 6]{Collins1996} and Shkarin \cite[Remarks (2)]{Shkarin2004}).

The paper is organized as follows. Section \ref{sec:defs} introduces the basic definitions and notation. Section \ref{sec:cfgos} studies cardinal invariants of GO spaces and introduces the \wob/-point technique. Section \ref{sec:hcomp} proves the main Theorem \ref{t:main:1} on compact subsets of homogeneous GO spaces, together with the results on semitopological GO groups. Sections \ref{sec:lex}--\ref{sec:blexhome} are devoted to constructions based on the lexicographic order on products of linearly ordered sets and on the space of binary sequences, which are later used to build examples. Section \ref{sec:exos} constructs homogeneous extensions of LOTS, and Section \ref{sec:bcs} uses them to build the space $\Hs$ with the large compact subset $\sS$. Section \ref{sec:qe} contains examples illustrating the sharpness of the main bounds, together with open questions.

%%% end file('ru//sec/intro.tex') %%%

\section{Terminology and definitions} \label{sec:defs}
%%% begin file('ru//sec/defs.tex') %%%

\subsection{Linearly ordered sets and spaces}

On a linearly ordered set $(X,<)$ the family of \term{open intervals} $\ilpp(a,b)=\set{x\in X: a<x<b}$, $\ilpp(-\infty,b)=\set{x\in X: x<b}$ and $\ilpp(a,+\infty)=\set{x\in X: a<x}$ for $a,b\in X$ is a base for a topology $\tp_{<}$. We say that the topology $\tp_{<}$ is generated by the linear order $<$.
We shall also use notation analogous to that for real numbers, e.g. $\ilcp[a,b)=\set{x\in X: a\leq x <b}$, and similarly for other interval types.

A subset $M\subset X$ is called \term{convex} if $a,b\in M$ and $a<b$ imply $\ilcc[a,b]\subset M$.

A topological space $X$ is called a \term{linearly ordered topological space} (LOTS) if there is a linear order on $X$ that generates the topology of $X$.
Subspaces of a LOTS are called \term{generalized ordered} (GO) topological spaces.

In this paper we shall also use LOTS to mean a linearly ordered set $(X,<)$ equipped with the topology $\tp_<$.

A LOTS $X$ is called \term{bounded} if $X$ has a minimum and a maximum.

% On a linearly ordered set $(X,<)$ a base of the topology is given by the intervals $\ilpp(a,b)=\set{x\in X: a<x<b}$, $\ilpp(-\infty,b)=\set{x\in X: x<b}$ and $\ilpp(a,+\infty)=\set{x\in X: a<x}$ for $a,b\in X$. Such spaces are called \term{linearly ordered topological spaces} (LOTS). Subspaces of a LOTS are called \term{generalized ordered} (GO) topological spaces.

A map $f: X\to Y$ between linearly ordered sets is called \term{increasing} (\term{decreasing}) if $f(a)\leq f(b)$ ($f(a)\geq f(b)$) for $a,b\in X$ with $a<b$. A map $f$ is called \term{strictly increasing} (\term{strictly decreasing}) if $f(a)< f(b)$ ($f(a)> f(b)$) for $a,b\in X$ with $a<b$. A map $f$ is called \term{(strictly) monotone} if $f$ is a (strictly) increasing or (strictly) decreasing map. A map $f$ is called an \term{order isomorphism} if $f$ is an increasing bijection. A map $f$ is called an \term{order reversing isomorphism} if $f$ is a decreasing bijection.

Note that an order isomorphism and an order reversing isomorphism are homeomorphisms between the LOTS $X$ and $Y$.

For a linearly ordered set $(X,<)$ denote by $\rev X=(X,<_{\rev})$ the set $X$ with the reverse order, $x<_{\rev} y$ if $y<x$. An order reversing isomorphism $f: X\to Y$ is an order isomorphism of the LOTS $\rev X\to Y$ and $X\to \rev Y$.

By an \term{order automorphism} we mean an order isomorphism from $X$ to itself.
By an \term{order reversing automorphism} we mean an order reversing isomorphism from $X$ to itself.
% Denote 
Denote
\begin{align*}
\oip X &= \set{f: f\text{ is an order automorphism of }X},
\\
\oim X &= \set{f: f\text{ is an order reversing automorphism of }X},
\\
\oipm X &= \oip X \cup \oim X.
\end{align*}
If $\oim X\neq\es$, then $\oip X$ is a subgroup of index $2$ in $\oipm X$.
We call a LOTS $X$ \term{order-homogeneous} if $\oip X$ acts transitively on $X$, that is, for any $x,y\in X$ there is an order isomorphism $f: X\to X$ with $f(x)=y$.
We call a LOTS $X$ \term{$\pm$-order-homogeneous} if $\oipm X$ acts transitively on $X$, that is, for any $x,y\in X$ there is either an order isomorphism $f: X\to X$ or an order reversing isomorphism $f: X\to X$ with $f(x)=y$. Clearly,
\[
\text{order-homogeneity }\rarr \text{ $\pm$-order-homogeneity }\rarr \text{homogeneity.}
\]

Let $X$ and $Y$ be LOTS. On the product $X\times Y$ we consider the \term{lexicographic} order: $(a,b)<(c,d)$ if $a<b$, or $a=b$ and $c<d$. By the \term{sum} $X+Y$ of the LOTS $X$ and $Y$ we mean the disjoint union $X \sqcup Y$ with the order such that $X$ and $Y$ are order-embedded in $X+Y$ and $x<y$ whenever $x\in X$ and $y\in Y$.



% Denote by $\homp X$ the set of order isomorphisms 

A set $M\subset X$ is called \term{cofinal} if for every $x\in X$ there is $y\in M$ with $x\leq y$.
The \term{cofinality} $\cf X$ of a linearly ordered set $X$ is the least cardinality of a cofinal subset of $X$. If $X$ has a maximal element, then $\cf X=1$. If $X$ is the empty set, we set $\cf X=0$.
Note that $\cf X$ is a regular cardinal and there is an increasing cofinal sequence $\set{x_\al:\al < \cf X}\subset X$.
For $x\in X$ denote by $\cf_l(x,X)$ the cofinality of the set
%$\lL(x,X)=\set{y\in X: y< x}$
$\lL(x,X)$  
and by $\cf_r(x,X)$ the cofinality of the set
%$\lR(x,X)=\set{y\in X:  x<y}$
$\lR(x,X)$  
in the reverse order, $\cf_r(x,X) = \cf_l(x,\rev X)=\cf\rev \lR(x,X)$.

A pair $(A,B)$ of subsets of $X$ is called a \term{cut} of the linearly ordered set $X$ if $A\cup B=X$, $A\cap B=\es$, and $a<b$ for all $a\in A$ and $b\in B$. A cut $(A,B)$ is called \term{proper} if $A\neq\es$ and $B\neq\es$. A proper cut is called a \term{gap} if $A$ has no maximum and $B$ has no minimum. A proper cut is called a \term{jump} if $A$ has a maximum and $B$ has a minimum. A jump is uniquely determined by the pair of points $a=\max A<b=\min B$, for which the interval $\ilpp(a,b)$ is empty. We call the point $a$ the \term{left jump point} and the point $b$ the \term{right jump point}. A left or right jump point will be called a \term{jump point}. A point $a$ is a left jump point if and only if $\cf_r(a,X)=1$, and a point $b$ is a right jump point if and only if $\cf_l(a,X)=1$. A whole jump is recovered from its left and right points: 
%$a=\max \lL(b,X)$ and $b=\min \lR(b,X)$.
$A=\ilpc(-\infty,a]$ and $B=\ilcp[b,+\infty)$.
Denote 
\[
\jumps X = \set{(a,b)\in X\times X:  a<b,\   \ilpp(a,b)=\es}.
\]

\subsection{Topological spaces}

A space $X$ is called \term{monotonically normal} if for every point $x\in X$ and every neighbourhood $U$ of $x$ a neighbourhood $g(x,U)$ of $x$ is defined, in such a way that $g(x,U)\cap g(y,V)=\es$ whenever $x\notin V$ and $y\notin U$. Every GO space is monotonically normal \cite[Theorem 5.21]{gru1984}. We call a space $X$ \term{power homogeneous} if $X^\tau$ is homogeneous for some nonzero cardinal $\tau$. A family of open sets $\cU$ of a space $X$ is called an \term{outer base} of a set $F\subset X$ if $F\subset U$ for some $U\in\cU$, and for every neighbourhood $W\supset F$ of $F$ there is $U\in\cU$ with $U\subset W$.
A space $X$ has \term{pointwise countable type} if for every $x\in X$ there is a compact set $K\subset X$ with $x \in K$ that has a countable outer base. An outer base of a singleton $\sset x$ is called a base of the point $x$.
%A family of open sets $\cU$ of a space $X$ is called \term{non-archimedean} if the following condition holds: if $U,V\in\cU$ and $U\cap V\neq\es$, then either $U\subset V$ or $V\subset U$. Spaces admitting a non-archimedean base are called \term{non-archimedean spaces}.
For a cardinal $\tau$, a set $G\subset X$ is of \term{type $G_\tau$} if $G=\bigcap \cU$ for some family of open sets $\cU$ of cardinality at most $\tau$. Sets of type $G_\om$ are called sets of \term{type $G_\de$}.
Denote by $\D_X=\set{(x,x):x\in X}$ the diagonal of $X\times X$.

Let $x\in X$ be a non-isolated point of a space $X$.
The \term{character} $\chi(x,X)$ of the point $x$ is the least cardinality of a base of $x$.
A family $\cU$ of nonempty open subsets of $X$ is called a \term{$\pi$-base} of a point $x\in X$ if for every neighbourhood $V$ of $x$ there is $U\in\cU$ with $U\subset V$.
The \term{$\pi$-character} $\pi\chi(x,X)$ of the point $x$ is the least cardinality of a $\pi$-base of $x$.
The \term{pseudocharacter} $\psi(x,X)$ of the point $x$ is the least cardinality of a family $\cU$ of neighbourhoods of $x$ with $\sset x=\bigcap\cU$.
Set $P(x,X)$ \cite{JuhaszNyikosSzentmiklossy2005} to be the least cardinality of a family $\cU$ of neighbourhoods of $x$ with $x\notin\Int \bigcap\cU$.
%
Set $T(x,A)=\min\set{|M|: M\subset A,\ x\in\cl M}$ for $x\in \cl A\subset X$,
$\sp_t(x,X)=\set{T(x,A): A\subset X\text{ and }x\in \cl A \setminus A}$,
$t(x,X)=\sup (\sp_t(x,X)\cup\sset \om)$,  
$\sp_t(X)=\bigcup\set{\sp_t(x,X): x\in X}$. 
%
Define the \term{weak tightness} $wt(x,X)=T(x,X\setminus \sset x)$ \cite{WilliamsZhou1998}.

If $x$ is an isolated point, set $\chi(x,X)=\pi\chi(x,X)=\psi(x,X)=t(x,X)=P(x,X)=wt(x,X)=\om$ and $\sp_t(x,X)=\es$.

We define the following cardinal invariants of a space $X$:
the \term{character}
$\chi(X)=\sup \set{\chi(x,X): x\in X}$,
the \term{$\pi$-character}
$\pi\chi(X)=\sup \set{\pi\chi(x,X): x\in X}$,
the \term{pseudocharacter}
$\psi(X)=\sup \set{\psi(x,X): x\in X}$,
$P(X)=\sup\set{P(x,X): x\in X}$, 
the \term{tightness}
$t(X)=\sup\set{t(x,X): x\in X}$, 
the \term{weak tightness}
$wt(X)=\sup\set{wt(x,X): x\in X}$,
the \term{diagonal number}
$\D(X)=\min\{\,\tau:$ the diagonal $\D_X$ in $X\times X$ has type  $G_\tau\,\}$
of the space $X$.



A sequence $\set{x_\al:\al<\tau}\subset X$ is called a \term{free sequence} if $\cl{\set{x_\al:\al<\la}}\cap \cl{\set{x_\al:\la \leq \al<\tau}}=\es$ for $\la<\tau$.

We call a point $x$ a \term{\wob/-point} if $x$ has a neighbourhood base $\cU$ that is linearly ordered by inclusion, that is, if $A,B\in\cU$, then either $A\subset B$ or $B\subset A$.
A \wob/-point has a well-ordered base: from a linearly ordered base one may take a cofinal well-ordered subfamily $\set{U_\al:\al<\la}\subset \cU$, where $\la=\cf \cU$ and $U_\be\subset U_\al$ for $\al<\be<\ga$.


Let $\tau$ be a cardinal. A point $x$ of a space $X$ is called a \term{$P_\tau$-point} if the intersection $\bigcap\cU$ of any family $\cU$, $|\cU|<\tau$, of neighbourhoods of $x$ is a neighbourhood of $x$. A space $X$ is called a \term{$P_\tau$-space} if every point of $X$ is a $P_\tau$-point. Every space is a $P_\om$-space. A $P_{\om_1}$-point is called a \term{$P$-point}. A $P_{\om_1}$-space is called a \term{$P$-space}.

Let $G$ be a group with a topology. The group $G$ is called \term{right-topological} if all right translations $\rho_g: G\to G$, $h\mapsto hg$ are continuous. The group $G$ is called \term{semitopological} if multiplication in $G$ is separately continuous. The group $G$ is called \term{quasitopological} if multiplication is separately continuous and taking the inverse in $G$ is continuous. The group $G$ is called \term{paratopological} if multiplication in $G$ is continuous. The group $G$ is called \term{topological} if multiplication and taking the inverse in $G$ are continuous.








%%% end file('ru//sec/defs.tex') %%%

\section{Cardinal invariants of GO spaces} \label{sec:cfgos}
%%% begin file('ru//sec/cfgos.tex') %%%


\begin{proposition}\label{p:1}
Let $X$ be a Hausdorff space and $x\in X$.
\begin{enumerate}
\item
If $X$ is a homogeneous space, then $\sp_t(x,X)=\sp_t(X)$.
\item
If $x\in K\subset X$, then $\sp_t(x,K)\subset \sp_t(x,X)$ and $\sp_t(K)\subset \sp_t(X)$.
\item
If $X$ is an infinite compact space, then $\sp_t(X)$ contains all infinite regular cardinals not exceeding $t(X)$.
\item
If $X$ is an infinite compact space, then $\om\in \sp_t(X)$.
\item
If $x$ is a non-isolated point, then
\begin{align*}
wt(x,X)&=\min \sp_t(x,X),
&
t(x,X)&=\sup \sp_t(x,X).
\end{align*}
\item
$|\sp_t(x,X)|\leq 1$ if and only if $wt(x,X)=t(x,X)$.
\end{enumerate}
\end{proposition}
\begin{proof}
(1) and (2) are obvious.

(3) Let $\tau=t(X)$. Since $X$ is compact, there is a free sequence $\set{x_\al:\al<\tau}\subset X$ \cite[Theorem 7.12]{Hodel1984handbook}.
Let $\mu\leq\tau$ be an infinite regular cardinal and let $y$ be a complete accumulation point of the set $A=\set{x_\al:\al<\mu}$. Then $\mu=T(y,A)\in \sp_t(X)$.

(4) There is a countable $A\subset X$ and $y\in \cl A\setminus A$. Then $\om=T(y,A)\in \sp_t(X)$.

(5) is obvious. (6) follows from (5).
\end{proof}


The following diagram shows the relations between the pointwise cardinal invariants. An arrow $A\to B$ means that $A(x,X)\geq B(x,X)$ for every Hausdorff space $X$ and $x\in X$.
\[
\begin{tikzcd}
\chi \ar[r] \ar[dd] \ar[rd]	&	\psi \ar[r]	&	P 
\\
 & \pi\chi \ar[rd]
\\
t \ar[rr]	&&	wt \ar[uu]
\end{tikzcd}
\]


% Из определений вытекает
% \begin{align*}\end{align*}
% \begin{gather*} \end{gather*}



\begin{proposition}\label{p:1-1}
Let $X$ be a Hausdorff space, $x\in X$.
\begin{enumerate}
\item
If $x$ is a \wob/-point, then
\[
\chi(x,X)=\pi\chi(x,X)=\psi(x,X)=t(x,X)=wt(x,X)=P(x,X).
\]
\item
If $\chi(x,X)=P(x,X)$, then $x$ is a \wob/-point.
\end{enumerate}
\end{proposition}
\begin{proof}
(1) Let $\la=P(x,X)$. It suffices to show that $\chi(x,X)\leq \la$.
Let $\cU$ be a linearly ordered base of $x$ and let $\set{V_\al:\al<\la}$ be a system of neighbourhoods of $x$ with $x\notin \Int \bigcap_{\al<\la}V_\al$. Take $\cO=\set{U_\al:\al<\la}\subset \cU$ such that $U_\al\subset V_\al$ and $U_\al\supset U_\be$ for $\al<\be<\la$. Then $\cO$ is a base at $x$.

(2) Let $\la=P(x,X)$. Let $\set{W_\al:\al<\la}$ be a base of $x$ and let $\set{V_\al:\al<\la}$ be a system of neighbourhoods of $x$ with $x\notin \Int \bigcap_{\al<\la}V_\al$. 
Let $\cU=\set{U_\al:\al<\la}$ be a family of open subsets of $X$ such that $U_\al\subset W_\al\cap V_\al$ and $U_\al\supset U_\be$ for $\al<\be<\ga$. Then $\cU$ is a linear base at $x$.
\end{proof}



\begin{proposition}\label{p:1+0}
Let $Y$ be a LOTS and $x\in X\subset Y$.
On $X$ we consider the order inherited from $Y$ and the subspace topology of the space $Y$.
Denote $L_*=X\cap \lL(x,Y)$, $R_*=X\cap \lR(x,Y)$, $L=L_*\cup \sset x$, $R=R_*\cup \sset x$.
\begin{enumerate}
\item
The point $x$ in $L$ is a \wob/-point. If $x\in \cl{L_*}$, then $\chi(x,L)=\cf_l(x,X)=\cf_l(x,Y)$.
\item
The point $x$ in $R$ is a \wob/-point. If $x\in \cl{R_*}$, then $\chi(x,R)=\cf_r(x,X)=\cf_r(x,Y)$.
\item
The point $x$ in $X$ is a \wob/-point if and only if one of the following conditions holds:
\begin{enumerate}
%\item
%$x\notin \cl{L_*}$ and $x\notin \cl{R_*}$, that is, $x$ is an isolated point of $X$;
\item
$x\notin \cl{L_*}$;
\item
$x\notin \cl{R_*}$;
\item
$x\in \cl{L_*}\cap \cl{R_*}$ and $\cf_l(x,X)=\cf_r(x,X)$.
\end{enumerate}
\item
%For $x\in X$,
\begin{enumerate}
\item
$\sp_t(x,X)=\es$ if 
$x\notin \cl{L_*}$ and $x\notin \cl{R_*}$, that is, $x$ is an isolated point of $X$;
\item
$\sp_t(x,X)=\sset{\cf_l(x,X)}$ 
if $x\in \cl{L_*}$ and $x\notin \cl{R_*}$;
\item
$\sp_t(x,X)=\sset{\cf_r(x,X)}$ 
if $x\notin \cl{L_*}$ and $x\in \cl{R_*}$;
\item
$\sp_t(x,X)=\sset{\cf_l(x,X),\cf_r(x,X)}$ 
if $x\in \cl{L_*}$ and $x\in \cl{R_*}$.
\end{enumerate}
\end{enumerate}
\end{proposition}
\begin{proof}
Let us prove (1).
Let $\la=\cf_l(x,X)$ and let $\set{x_\al:\al<\la}$ be a set cofinal in $L_*$. Then
the family $\set{\ilpc(x_\al,x]:\al<\la}$ forms a base of $x$ in $\ilpc(-\infty, x]$, and
the family $\set{\ilpc(x_\al,x]\cap X:\al<\la}$ forms a base of $x$ in $L$.

Item (2) follows from (1).

Let us prove (3). ($\larr$) 
Item (a) follows from (2). Item (b) follows from (1).
Let us prove (c). Let $\la=\cf_l(x,X)$, let $\set{x_\al:\al<\ga}\subset L_*$ be a sequence increasing to $x$, and let $\set{y_\al:\al<\ga}\subset R_*$ be a sequence decreasing to $x$.
Then $\set{\ilpp(x_\al,y_\al)\cap X:\al<\ga}$ is a linearly ordered base at the point $x$.

($\rarr$) If $x\notin \cl{L_*}$ or $x\notin \cl{R_*}$, then either (a) or (b) holds.
Consider the case $x\in \cl{L_*}\cap \cl{R_*}$. There is a linearly ordered base $\set{\ilpp(x_\al,y_\la)\cap X:\al<\la}$ of $x$, where $\la$ is some regular cardinal and $x_\al\in L_*$, $y_\al\in R_*$ for $\al<\la$. Then $\la=\cf_l(x,X)=\cf_r(x,X)$, and (d) holds.

Item (4) is obvious.
\end{proof}

The following statement follows from Proposition \ref{p:1+0}(3).


\begin{proposition}\label{p:1+0.5}
Let $X$ be a LOTS and $x\in X$. If $x$ is a minimum, a maximum, or a jump point, then $x$ is a \wob/-point.
\end{proposition}



\begin{proposition}\label{p:1+1}
Let $X$ be a GO space and $x\in X$.
Then  
%$1\leq |\sp_t(x,X)|\leq 2$,
\begin{align*}
\pi\chi(x,X) &= wt(x,X)= P(x,X),
&
\chi(x,X) &= \psi(x,X) = t(x,X).
\end{align*}
If $x$ is not an isolated point of $X$, then
\begin{align*}
\min \sp_t(x,X) &= \pi\chi(x,X),
&
\max \sp_t(x,X) &= \chi(x,X).
\end{align*}

\end{proposition}
\begin{proof}
The proposition is obvious if $x$ is an isolated point. Consider the case when $x$ is not an isolated point. Then $\sp_t(x,X)\neq\es$.

If $|\sp_t(x,X)|=1$, then $\min \sp_t(x,X)=\max \sp_t(x,X)$, and Proposition \ref{p:1+0}(3) implies that $x$ is a \wob/-point. The required equality then follows from Proposition \ref{p:1-1} and Proposition \ref{p:1+0}(1,2).

Consider the case $|\sp_t(x,X)|>1$.
Let $Y$ be a LOTS in which $X$ is embedded.
On $X$ we consider the order inherited from $Y$.
Denote $L_*=X\cap \lL(x,Y)$, $R_*=X\cap \lR(x,Y)$, $L=L_*\cup \sset x$, $R=R_*\cup \sset x$, $\la=\cf_l(x,X)$, $\rho=\cf_r(x,X)$.
By Proposition \ref{p:1+0}(4), $x\in \cl{L_*}\cap \cl{R_*}$, $|\sp_t(x,X)|=2$, and $\sp_t(x,X)=\sset{\la,\rho}$. Without loss of generality $\la<\rho$. Then $\la=\min \sp_t(x,X)$ and $\rho=\max \sp_t(x,X)$. By Proposition \ref{p:1+0}(1,2), $x$ is a \wob/-point in $L$ and in $R$.
By Proposition \ref{p:1-1}, $\pi\chi(x,L) = wt(x,L)= P(x,L)=\la$ and $\pi\chi(x,R) = wt(x,R)= P(x,R)=\rho>\la$. Hence $\pi\chi(x,X) = wt(x,X)= P(x,X)=\la$.

By Proposition \ref{p:1-1}, $\chi(x,L) = t(x,L)= \psi(x,L)=\la<\rho$ and $\chi(x,R) = t(x,R)= \psi(x,R)=\rho$. Hence $\chi(x,X)= \psi(x,X) = t(x,X) = \rho$.
\end{proof}

The following statement follows from Propositions \ref{p:1+0} and \ref{p:1+1}, and is also easily verified directly.

\begin{proposition}\label{p:1+1.5}
Let $X$ be a LOTS and $x\in X$. Then
\begin{align*}
\pi\chi(x,X) &= \om + \min \cf_l(x,X), \cf_r(x,X),
\\
\chi(x,X) &=  \om + \max \cf_l(x,X), \cf_r(x,X).
\end{align*}
\end{proposition}



The following statement follows from Propositions \ref{p:1-1}, \ref{p:1+0} and \ref{p:1+1}.

\begin{proposition}\label{p:1+2}
Let $X$ be a GO space and $x\in X$.
The following conditions are equivalent:
\begin{enumerate}
\item
$x$ is a \wob/-point;
\item
$\chi(x,X)= \pi\chi(x,X)$;
\item
$|\sp_t(x,X)|\leq 1$.
\end{enumerate}
\end{proposition}

\begin{proposition}[{\cite[Theorem 2.5]{WilliamsZhou1998}, \cite[Theorem 3.12 (iii)]{Juhasz1992}}]\label{p:1+3}
In a monotonically normal compact space the set of \wob/-points is dense.
\end{proposition}

We give a simple proof of a corollary of this fact.

\begin{cor}\label{c:cfgos:1}
In a compact LOTS the set of \wob/-points is dense.
\end{cor}
\begin{proof}
Let $X$ be a compact LOTS. Let $U$ be a nonempty open subset of $X$.
Isolated points of $U$ are \wob/-points. Consider the case when $U$ has no isolated points.
Set $\cI=\set{\ilpp(a,b): a,b\in U,\ \es\neq\ilpp(a,b)\subset U}$. On $\cI$ consider the order: for $A,B\in \cI$, $A<B$ if $\cl B\subset A$.
By the Hausdorff maximum principle there is a maximal linearly ordered subset $\cU\subset \cI$. Let $F=\bigcap\cU$. Then $F\neq\es$. Let $a=\min F$ and $b=\max F$.

Let us show that $\ilpp(a,b)=\es$. Suppose otherwise. Since $U$ has no isolated points, the set $\ilpp(a,b)$ is infinite. Let $a',c,b'\in \ilpp(a,b)$, $a'<c<b'$. Then $I=\ilpp(a',b')\in \cI$ and $A<I$ for every $A\in \cU$, contradicting the maximality of $\cU$. Thus $\ilpp(a,b)=\es$.

If $a=b$, then $F=\sset{a}$ and $\cU$ is a linearly ordered base of the point $a$. If $a\neq b$, then $a$ and $b$ are jump points. By Proposition \ref{p:1+0.5}, $a$ and $b$ are \wob/-points.
\end{proof}

%The following statement is easily proved.

\begin{proposition}\label{p:1+0a}
Let $Y$ be a LOTS, $X\subset Y$.
On $X$ we consider the order inherited from $Y$ and the subspace topology of the space $Y$.
Let $L_x=\ilpp(-\infty,x)\cap X$ and $R_x=\ilpp(x,+\infty)\cap X$ for $x\in X$.
\begin{enumerate}
\item
The topology of $X$ coincides with the order topology on $X$ if and only if
for every $x\in X$ the following conditions hold:
\begin{enumerate}
\item
$\max L_x$ exists if $x\notin\cl{L_x}$;
\item
$\min R_x$ exists if $x\notin\cl{R_x}$.
\end{enumerate}
\item
The topology of $X$ coincides with the order topology on $X$ if
one of the following conditions holds:
\begin{enumerate}
\item
$x\in \cl{L_x}\cap \cl{R_x}$ for every $x\in X$;
\item
$X$ is dense in $Y$ and contains no jump points of $Y$;
\item
$X$ is compact.
\end{enumerate}
\end{enumerate}
\end{proposition}
\begin{proof}
Item (1) is verified directly without difficulty. Item (2) follows from (1).
\end{proof}

The following statements are well known \cite{Nagata1985}.

\begin{proposition}\label{p:1+4}
Let $X$ be a LOTS.
\begin{enumerate}
\item
The space $X$ is locally compact if $X$ has no gaps.
\item
The space $X$ is compact if and only if $X$ has no gaps and $X$ has a minimum and a maximum.
\item
The space $X$ is connected if and only if $X$ has no gaps and no jumps.
\end{enumerate}
\end{proposition}


\begin{proposition}\label{p:1+5}
Let $X$ be a compact LOTS.
The space $X$ is zero-dimensional if and only if for any $a,b\in X$, $a<b$, there is a jump $(A,B)$ such that $\ilpc(-\infty,a]\subset A$ and $\ilcp[b,+\infty)\subset B$.
\end{proposition}
\begin{proof}
($\rarr$) Let $a<b$. Since $X$ is a zero-dimensional compact space, there is a clopen $U\subset X$ such that $\ilpc(-\infty,a]\subset U$ and $\ilcp[b,+\infty)\subset X\setminus A$. Set $a'=\max U$. Then the set $\ilpc(-\infty,a']$ is clopen.
Set $b'=\min \ilpp(a',+\infty)$, $A=\ilpc(-\infty,a']$ and $B=\ilcp[b,+\infty)$.
Then $(A,B)$ is a jump and $\ilpc(-\infty,a]\subset A$ and $\ilcp[b,+\infty)\subset B$.

($\larr$) The condition implies that $X$ is totally disconnected. Since $X$ is compact, $X$ is zero-dimensional \cite[Theorem 6.2.9]{EngelkingBookRu}.
\end{proof}



\begin{proposition}\label{p:1+6}
Let $X$ be a zero-dimensional compact LOTS. A base at a point $x\in X$ is formed by open intervals of the form $\ilcc[a,b]$, where $a$ is the minimum of $X$ or a right jump point, and $b$ is the maximum of $X$ or a left jump point.
\end{proposition}
\begin{proof}
Let $U$ be an open interval containing $x$.

Let us find the point $a$.
If $\ilpp(-\infty,x)\cap U=\es$, set $a=x$. Otherwise, take $a'\in \ilpp(-\infty,x)\cap U$. By Proposition \ref{p:1+5} there exist $a'',a\in X$ such that $a'\leq a''<a\leq x$ and $\ilpp(a'',a)=\es$.

Let us find the point $b$.
If $\ilpp(x,+\infty)\cap U=\es$, set $b=x$. Otherwise, take $b'\in \ilpp(x,+\infty)\cap U$. By Proposition \ref{p:1+5} there exist $b'',b\in X$ such that $x\leq b<b''\leq b'$ and $\ilpp(b,b'')=\es$.
\end{proof}








%%% end file('ru//sec/cfgos.tex') %%%

\section{Compact subsets of homogeneous GO spaces} \label{sec:hcomp}
%%% begin file('ru//sec/hcomp.tex') %%%

A locally compact power homogeneous monotonically normal space has countable character \cite[Corollary 3.22]{Arhangelskii2004}.
A power homogeneous GO space of pointwise countable type has countable character \cite[Corollary 29]{Arhangelskii2005}.

\begin{theorem}\label{t:main:1}
Let $X$ be a homogeneous GO space. Then
\begin{enumerate}
\item
$\chi(K)\leq\om_1$ and $|K|\leq 2^{\om_1}$ for every compact $K\subset X$;
\item
if there is a compact $K\subset X$ with $\chi(K)>\om$, then $\chi(X)=\om_1$ and $\pi\chi(X)=\om$.
\end{enumerate}
\end{theorem}
\begin{proof}
Let $K\subset X$ be compact.
By Proposition \ref{p:1+0}(4) and Proposition \ref{p:1}(1), $|\sp_t(X)|\leq 2$.
By Proposition \ref{p:1}(2), $\sp_t(K)\subset \sp_t(X)$, hence $|\sp_t(K)|\leq 2$.
By Proposition \ref{p:1+0}(3), $\sp_t(K)\subset\sset{\om,\om_1}$.

Let us prove (1). If $K$ is finite, the statement is obvious. Consider the case when $K$ is infinite.
By Proposition \ref{p:1+1}, $\chi(K)=\max \sp_t(K)\leq \om_1$.
By the Arhangel'skii theorem
\cite[Theorem 4.5]{Hodel1984handbook}, $|K|\leq 2^{\om_1}$.

Let us prove (2). Let $\chi(K)>\om$. By (1), $\chi(K)=\om_1\in \sp_t(K)$.
By Proposition \ref{p:1}(4), $\om\in \sp_t(K)$. Since $\sp_t(K)\subset\sset{\om,\om_1}$, we get $\sp_t(K)=\sset{\om,\om_1}$. Since $\sp_t(K)\subset \sp_t(X)$ and $|\sp_t(X)|\leq 2$, we get $\sp_t(X)=\sset{\om,\om_1}$. By Proposition \ref{p:1+1}, $\chi(X)=\om_1$ and $\pi\chi(X)=\om$.
\end{proof}


\begin{proposition}\label{p:main:1}
Let $\la$ be a regular cardinal and $C\subset \la$ a stationary subset. Then $\D(C)=\la$.
\end{proposition}
\begin{proof}
Since $|C|=\la$, we have $\D(C)\leq\la$. Suppose $\D(C)\neq\la$. Then there is a family $\cW$ of open neighbourhoods of the diagonal $\D_C$ in the square $C\times C$ such that $\bigcap\cW=\D_C$ and $|\cW|<\la$.

For $W\in\cW$ find $x_W\in C$ such that $U_W\times U_W\subset W$, where $U_W=\set{x\in C: x> x_W}$. Let $c_0=\min C$ and $C_*=C\setminus \sset{c_0}$. For each $x\in C_*$ there is $f(x)\in C$, $f(x)<x$, such that $V_x \times V_x\subset W$, where $V_x=\set{y\in C: f(x)< y \leq x}$. By Fodor's lemma, the set $f^{-1}(x_W)$ is stationary in $\la$ for some $x_W\in C$. Then $U_W\times U_W=\bigcup_{x\in f^{-1}(x_W)}V_x\times V_x\subset W$.

Since $\la$ is a regular cardinal and $|\cW|<\la$, there is $a\in C$ such that $x_W<a$ for all $W\in \cW$. Take $b\in C$, $a<b$. Then $(a,b)\notin \D_C$ and $(a,b)\in W$ for all $W\in\cW$, contradicting $\bigcap\cW=\D_C$.
\end{proof}


\begin{theorem}\label{t:main:1+1}
Let $G$ be a semitopological group that is a GO space.
Then $G$ is hereditarily paracompact and
\[
P(X) = wt(X)= \pi\chi(X) = \D(X) =
\chi(X) = \psi(X) = t(X).
\]
\end{theorem}
\begin{proof}
Let $\tau=P(X)$. By Proposition \ref{p:1+1} and the homogeneity of $X$, $wt(X)= \pi\chi(X) = \tau$.
Since $G$ is a semitopological group and $\pi\chi(X) = \tau$, we get $\D(X)\leq\tau$ (\cite[Theorem 1]{Reznichenko2023-1}, \cite[Theorem 2.4]{arh-rezn2005}, and \cite[Corollary 5.7.5]{at2009} for countable $\tau$). Since $\psi(X)\leq \D(X)$, Proposition \ref{p:1+1} and the homogeneity of $X$ give $\chi(X) = t(X) = \psi(X) \leq \D(X)\leq \tau$. Since $\pi\chi(X)\leq \chi(X)$, we get $\chi(X) = t(X) = \psi(X) = \D(X) = \tau$.

Let us prove that $G$ is hereditarily paracompact. Suppose otherwise. Since a GO space is monotonically normal, $G$ is a monotonically normal, non-hereditarily-paracompact space. By \cite[Theorem I]{BaloghRudin1992}, there is a regular cardinal $\la$ and a stationary subset $C\subset \la$ homeomorphic to some $F\subset G$. By Proposition \ref{p:main:1}, $\D(F)=\la$. Since $\D(F)\leq \D(X)$, we get $\la\leq \tau$. A stationary subset is not discrete, so let $x\in F$ be a non-isolated point of $F$. Then $\tau=wt(x,X)\leq T(x,F\setminus\sset x)<\la$, contradicting $\la\leq \tau$.
\end{proof}


\begin{theorem}\label{t:main:2}
Let $G$ be a semitopological group that is a GO space and not a $P$-space.
Then $G$ has countable character, is submetrizable, and is hereditarily paracompact.
\end{theorem}
\begin{proof}
Since $G$ is not a $P$-space, $P(x,G)=\om$ for some $x\in G$.
By Theorem \ref{t:main:1+1}, $\chi(G)=\D(G)=\om$ and $G$ is hereditarily paracompact.
Paracompact spaces with a $G_\de$-diagonal are submetrizable \cite[Corollary 2.9]{gru1984}.
Note also that GO spaces with a $G_\de$-diagonal are (hereditarily) paracompact \cite[Theorem 4.5]{Lutzer1971}.
\end{proof}

In $P$-spaces compact subsets are finite.
In submetrizable spaces compact subsets are metrizable.
Hence the following statement follows from Theorem \ref{t:main:2}.

\begin{cor}\label{с:main:1}
In semitopological GO groups compact subsets are metrizable.
\end{cor}

\begin{theorem}\label{t:main:3}
Let $G$ be a semitopological group that is a LOTS and not a $P$-space.
Then $G$ is metrizable.
\end{theorem}
\begin{proof}
By Theorem \ref{t:main:2}, $G$ has a $G_\de$-diagonal.
A LOTS with a $G_\de$-diagonal is metrizable \cite{Lutzer1969}, \cite[Theorem 2.3]{gru1984}.
Hence $G$ is metrizable.
\end{proof}


\begin{theorem}[{\cite[Theorem 2.6]{BuzyakovaVural2014}, \cite[Theorem 3.5]{WilliamsZhou1998}}]
\label{t:main:3+1}
Let $G$ be a monotonically normal paratopological group. Then $G$ is hereditarily paracompact and every compact subset of $G$ has countable tightness.
\end{theorem}
\begin{proof}
Monotonically normal paratopological groups are hereditarily paracompact \cite[Theorem 2.6]{BuzyakovaVural2014}, and hence compact subsets of such groups have countable tightness \cite[Theorem 3.5]{WilliamsZhou1998}.
\end{proof}



%\begin{theorem}\label{t:main:3+2} Compact subsets of monotonically normal topological groups are metrizable. \end{theorem} \begin{proof} Let $G$ be a monotonically normal topological group and $K\subset G$ compact. Let $H$ be the group generated by $K$. Then $H$ is a $\si$-compact topological group and, by Tkachenko's theorem \cite[Corollary 5.3.22]{at2009}, has countable Souslin number. A monotonically normal space with countable Souslin number is hereditarily Lindel\"of \cite[Theorem 3.1]{Moody1992}, \cite{Ostaszewski1978}. Hence $H$ is hereditarily Lindel\"of and therefore has countable pseudocharacter. Topological groups with countable pseudocharacter are submetrizable \cite[Theorem 3.3.16]{at2009}. Then $H$ is submetrizable and every compact subset of $H$ is metrizable \cite[Corollary 3.3.17]{at2009}. Hence $K$ is metrizable. \end{proof} 

\begin{theorem}\label{t:main:2+1}
Let $G$ be a topological group that is a GO space and not a $P$-space.
Then $G$ is metrizable.
\end{theorem}
\begin{proof}
By Theorem \ref{t:main:2}, $G$ has countable character. Since $G$ has countable character, the Birkhoff--Kakutani theorem \cite[Theorem 3.3.12]{at2009} implies that the group $G$ is metrizable.
\end{proof}

%2.3. THEOREM (LUTZER [1969]). Every linearly ordered space with a Gs-diagonal is metrizable.




%%% end file('ru//sec/hcomp.tex') %%%

%\section{Order-homogeneous and $\pm$-order-homogeneous LOTS} \label{sec:cssoh}
%%% begin file('ru//sec/cssoh.tex') %%%


\begin{theorem}\label{t:cssoh:1}
Let $X$ be an order-homogeneous LOTS.
Then every compact subset of $X$ is first countable.
\end{theorem}
\begin{proof}
Let $K\subset X$ be compact. Suppose $\chi(K)>\om$. Then $\om_1+1$ embeds monotonically in $K$. Without loss of generality, let $f: \om_1+1 \to K$ be an increasing embedding. Then $\om=\cf_l(f(\om),X)\neq \cf_l(f(\om_1),X)=\om_1$, contradicting the order-homogeneity of $X$.
\end{proof}


%%% end file('ru//sec/cssoh.tex') %%%

\section{The product of linearly ordered spaces with the lexicographic order} \label{sec:lex}
%%% begin file('ru//sec/lex.tex') %%%

%Let $X$ be a linearly ordered set, $\la$ an ordinal. Denote by $\lex \la X$ the set $X^\la$ with the lexicographic order: $f<g$ for distinct $f,g\in X^\la$ if $f(\al)<g(\al)$, where $\al=\min \set{\be<\la: f(\be)\neq g(\be)}$. Denote $\blex f\al\la X=\set{g\in X^\la: g(\be)=f(\be)\text{ for }\be<\al}$ for $\al<\ga=\dom f$.

Let $X$ be a linearly ordered set, $\la$ an ordinal. On the set $X^\la$ we consider the lexicographic order: $f<g$ for distinct $f,g\in X^\la$ if $f(\al)<g(\al)$, where $\al=\min \set{\be<\la: f(\be)\neq g(\be)}$. 
Denote $\blex f\al\la X=\set{g\in X^\la: g(\be)=f(\be)\text{ for }\be<\al}$ for $\al<\la=\dom f$.

On $X^\la$ we consider the order topology.

\begin{proposition}\label{p:2}
If $X$ is a linearly ordered set with no greatest and no least element, $\la$ is a limit ordinal, and $f\in\lex \la X$, then the family $\set{\blex f\al \la X: \al<\la}$ forms a base of open neighbourhoods of the point $f$ in the linearly ordered space $\lex \la X$.
\end{proposition}
\begin{proof}
Let us show that the set $\blex f\al \la X$ is open. Let $g\in \blex f\al \la X$. There exist $a,b\in \blex f\al \la X$ such that $a(\al)<g(\al)<b(\al)$. Then $g\in (a,b)\subset \blex f\al \la X$.

Let us show that the family $\set{\blex f\al \la X: \al<\la}$ forms a base of neighbourhoods of $f$. Let $a,b\in \lex \la X$ and $a< f <b$. Let $\al=\min\set{\de<\la: a(\de)\neq f(\de)}$, $\be=\min\set{\de<\la: b(\de)\neq f(\de)}$, $\te=\max \al,\be$. Then $f\in \blex f{\te+1} \la X\subset (a,b)$.
\end{proof}

\begin{proposition}\label{p:h:tg1}
%\begin{theorem}\label{t:h:1}
Let $\la$ be a regular cardinal and $G=\lex \la \R$ with the order topology.
Then $G$ is a LOTS topological group and $\chi(G)=\la$.
The space $G$ is a $P_\la$-space.
\end{proposition}
\begin{proof}
On $G$ we consider the group operation induced from $\R^\la$, addition of functions. By Proposition \ref{p:2}, the family of open subgroups $\set{\blex 0\al \la \R: \al<\la}$ forms a base of neighbourhoods of the zero function.
\end{proof}

The following statement follows from \cite[Proposition 1.3(b)]{AkinHrbacek2002}.

\begin{proposition}\label{p:h:1}
Let $\la$ be an ordinal and $X$ a compact LOTS.
Then the LOTS $\lex \la X$ is compact.
\end{proposition}


%%% end file('ru//sec/lex.tex') %%%

\section{Binary sequences with the lexicographic order} \label{sec:blex}
%%% begin file('ru//sec/blex.tex') %%%

%Denote by $\B$ the set $\sset{0,1}$ with the natural order and the discrete topology.
Henceforth we regard $\B$ as the linearly ordered space $\sset{0,1}$.
Let $\la$ be an ordinal.
Henceforth on $\lexab$ we consider the lexicographic order and the order topology.

Denote by $R_1$ the order-reversing bijection of $\B$ onto itself: $R_1(0)=1$ and $R_1(1)=0$.
Denote
\[
R_\la=(R_1)^\la: \lexab\to \lexab,\ f\mapsto R_1\circ f.
\]
The map $R_\la$ is
%a bijection of $\lexab$ onto itself, reversing the order.
an order reversing isomorphism of $\lexab$ onto itself.

Denote by $0_\la$ (resp., $1_\la$) the function on $\la$ identically equal to zero (resp., one). Note that $0_\la =\min \lexab$ and $1_\la =\max \lexab$.



\begin{proposition}\label{p:blex:1}
Let $\la$ be a regular cardinal, $x\in \lexab$, $\ka_0=\cf x^{-1}(0)$ and $\ka_1=\cf x^{-1}(1)$. Then 
\begin{enumerate}
\item
$\cf_l(x,\lexab)=\ka_1$, $\cf_r(x,\lexab)=\ka_0$, and $\la\in\sset{\ka_0,\ka_1}$;
\item
$\chi(x,\lexab)=\la$, 
\[
\pi\chi(x,\lexab)=\begin{cases}
\min \sset{\ka_0, \ka_1}, & \ka_0\text{ and }\ka_1\text{  are infinite,}
\\
\la, & \ka_0\text{ or }\ka_1\text{ is finite.}
\end{cases}
\]
\item
the point $x$ is a \wob/-point if and only if one of the following conditions holds:
\begin{enumerate}
\item
$|x^{-1}(0)|=|x^{-1}(1)|=\la$;
\item
$x^{-1}(0) =\es$  or $x^{-1}(1) =\es$;
\item
the set $x^{-1}(0)$ has a maximum;
\item
the set $x^{-1}(1)$ has a maximum.
\end{enumerate}
\item
the point $x$ is a left jump point if and only if the set $x^{-1}(0)$ has a maximum; denoting $\mu = \max x^{-1}(0)$, in this case $x=\restr x\mu \symblecnc 0 \symblecnc 1_\la$ and $r=\restr x\mu \symblecnc 1 \symblecnc 0_\la$ is the right jump point;
\item
the point $x$ is a right jump point if and only if the set $x^{-1}(1)$ has a maximum; denoting $\mu = \max x^{-1}(0)$, in this case $x=\restr x\mu \symblecnc 1 \symblecnc 0_\la$ and $l=\restr x\mu \symblecnc 0 \symblecnc 1_\la$ is the left jump point.
\end{enumerate}
\end{proposition}
\begin{proof}
Let us prove (1). Since $\la$ is a regular cardinal and $\la=x^{-1}(0)\cup x^{-1}(1)$, either $\ka_0=\la$ or $\ka_1=\la$.
Denote $M=x^{-1}(1)$. Then $\ka_1=\cf M$.

If $\ka_1=0$, then $x=0_\la$ and $\cf_l(x,\lexab)=\ka_1=0$.

Consider the case $\ka_1=1$. Then $\mu=\max M$ exists.
Set $l=\restr x\mu \symblecnc 0 \symblecnc 1_\la$.
Then $l$ is a left jump point, $x$ is a right jump point, $l=\max \ilpp(-\infty,x)$, and $\cf_l(x,\lexab)=\ka_1=1$.

Consider the case $\ka_1\geq\om$.
Let $\set{\eta_\al:\al<\ka_1}\subset M$ be a strictly increasing cofinal sequence.
For $\al<\ka_1$ define $y_\al\in \lexab$:
\begin{equation*}
y_\al(\nu) = 
\begin{cases}
x(\nu),	& \nu\neq \eta_\al; 
\\
0,	& \nu= \eta_\al.
\end{cases}
\end{equation*}
Then the strictly increasing sequence $\set{y_\al:\al<\ka}\subset \lL(x,\lexab)$ converges to $x$.
Hence $\cf_l(x,\lexab)=\ka_1$.

Since $x^{-1}(0)=\left(R_\la(x)\right)^{-1}(1)$, we have
\[
\cf_r(x,\lexab)=\cf_l(R_\la(x),\lexab)=\cf \left(R_\la(x)\right)^{-1}(1)= \cf x^{-1}(0)=\ka_0.
\]

Let us prove (2). By (1) and Propositions \ref{p:1+0} and \ref{p:1+1}, $\chi(x,\lexab)=\la$. If $\ka_0$ or $\ka_1$ is finite, then $\sp_t(\lexab)=\sset \la$, and by Proposition \ref{p:1+1}, $\pi\chi(x,\lexab)=\la$. Otherwise, $\sp_t(\lexab)=\sset{\ka_0,\ka_1}$, and by Proposition \ref{p:1+1}, $\pi\chi(x,\lexab)=\min \ka_0,\ka_1$.

Let us prove (3). By (2), each of conditions (a), (b), (c), (d) is equivalent to $\pi\chi(x,\lexab)=\chi(x,\lexab)$. By Proposition \ref{p:1+2}, the latter equality is equivalent to $x$ being a \wob/-point.

Let us prove (4). The point $x$ is a left jump point if and only if $\cf_r(x,\lexab)=1$, that is, by (1), $\cf x^{-1}(0)=\ka_0=1$. The condition $\cf x^{-1}(0)=1$ is equivalent to $x^{-1}(0)$ having a maximal element. Let $\mu = \max x^{-1}(0)$. Then $x=\restr x\mu \symblecnc 0 \symblecnc 1_\la$, and the right jump point $r=\min \lR(x,X)$ equals $\restr x\mu \symblecnc 1 \symblecnc 0_\la$.

Item (5) follows from (4).
\end{proof}

For a regular cardinal $\la$ and a cardinal $\ka$ set
\begin{align*}
\homg l,\ka\la &= \set{ x\in \lexab: \cf_l(x,\lexab) = \ka \text{ and } \cf_r(x,\lexab) = \la},
\\
\homg r,\ka\la &= \set{ x\in \lexab: \cf_r(x,\lexab) = \ka  \text{ and } \cf_l(x,\lexab) = \la},
\\
\homg lr,\ka\la &= \homg l,\ka\la \cup \homg r,\ka\la,
\\
\homm\la &= \set{ x\in \lexab: \cf_r(x,\lexab) = \la  \text{ and } \cf_l(x,\lexab) = \la},
\end{align*}

Clearly, $\homm\la = \homg l,\la\la = \homg r,\la\la = \homg lr,\la\la$.

Note that by Proposition \ref{p:blex:1}(1), if $\ka$ is a regular cardinal and $\ka<\la$, then $\homg l,\ka\la = \set{ x\in \lexab: \cf_l(x,\lexab) = \ka}$ and $\homg r,\ka\la = \set{ x\in \lexab: \cf_r(x,\lexab) = \ka}$.


\begin{proposition}\label{p:blex:1+1}
Let $\la$ be a regular cardinal.
\begin{enumerate}
\item
$\homg l,0\la=\sset{1_\la}$ and $\homg r,0\la=\sset{0_\la}$;
\item
$\homg l,1\la$ is the set of right jump points and $\homg r,1\la$ is the set of left jump points; these sets are dense in $\lexab$;
\item
if $\ka\leq \la$ is a regular cardinal, then 
\begin{enumerate}
\item
the sets $\homg l,\ka\la$ and $\homg r,\ka\la$ are dense in $\lexab$;
\item
the topologies of $\homg l,\ka\la$, $\homg r,\ka\la$ and $\homg lr,\ka\la$ coincide with the order topology with respect to the order inherited from $\lexab$; accordingly, the spaces listed are LOTS;
\item
if $X\in\sset{\homg l,\ka\la,\homg r,\ka\la,\homg lr,\ka\la}$ and $x\in X$, then $\chi(x,X)=\la$ and $\pi\chi(x,X)=\ka$;
\item
the spaces $\homg l,\ka\la$, $\homg r,\ka\la$ and $\homg lr,\ka\la$ are $P_\ka$-spaces;
\end{enumerate}
\item
if $\ka< \la$ is a regular cardinal, then 
\begin{enumerate}
\item
the sets $\homg l,\ka\la$ and $\homg r,\ka\la$ are disjoint;
\item
$\homg lr,\ka\la=\set{x\in\lexab: \pi\chi(x,\lexab)=\ka}$.
\end{enumerate}
\item
the family
\begin{gather*}
\{ \homg l,0\la, \homg r,0\la, \homg l,1\la, \homg r,1\la, \homm\la \} \cup 
\\
\set{\homg l,\ka\la, \homg r,\ka\la: \ka<\la\text{ is a regular cardinal}}
\end{gather*}
is a partition of $\lexab$;
\item
if $\ka>\om$ is a regular ordinal, then compact subsets of $\homg l,\ka\la$, $\homg r,\ka\la$ and $\homg lr,\ka\la$ are finite;
\item
the space $\homg lr,\om\la$ is countably compact;
\item
the transfinite line $\om_1+1$ embeds in the space $\homg lr,\om{\om_1}$.
\end{enumerate}
\end{proposition}
\begin{proof}
Items (1) and (2) follow from Proposition \ref{p:blex:1} (1) and (2).


Let us prove (3). Let us prove (3)(a). Let $U$ be a nonempty open subset of $\lexab$. By Proposition \ref{p:1+6}, there exist $a,b\in \lexab$ such that $\ilcc[a,b]\subset U$ is a nonempty open subset of $\lexab$, $a$ is the minimum of $\lexab$ or a right jump point, and $b$ is the maximum of $\lexab$ or a left jump point. By Proposition \ref{p:blex:1}(1), $\cf_r(a,\lexab)=\cf_l(b,\lexab)=\la$. Let $\sq{x_\al}{\al<\la}\subset \ilcc[a,b]$ be a strictly increasing sequence converging to the point $b$. Let $y$ be the limit of the sequence $\sq{x_\al}{\al<\ka}$. Then $y\in \ilcc[a,b] \cap \homg l,\ka\la\subset U$. Similarly, $\homg r,\ka\la\cap U\neq\es$.

Let us prove (3)(b). The sets $\homg l,\ka\la$ and $\homg r,\ka\la$ are dense in $\lexab$ and contain no jump points and no minimum or maximum. Then (b) follows from Proposition \ref{p:1+0a}.

Item (3)(c) follows from Proposition \ref{p:blex:1} (1) and (2).

Let us prove (3)(d). By Proposition \ref{p:1+1} and item (4)(c), $P(x,\lexab)=\ka$ for all $x\in \homg lr,\ka\la$. Hence $\homg lr,\ka\la$ is a $P_\ka$-space.


Let us prove (4). Item (4)(a) follows from $\ka<\la$ and Proposition \ref{p:blex:1} (1) and (2).
Item (4)(b) follows from $\ka<\la$, Proposition \ref{p:blex:1} (1), (2), and Proposition \ref{p:1+1}.

Item (5) follows from the previous items.

Let us prove (6). By (3)(d), $\lexab$ is a $P$-space. Hence compact subsets of $\homg lr,\ka\la$ are finite.

Let us prove (7). Let $M\subset \homg lr,\om\la$ and $|M|=\om$. Let $F$ be the closure of $M$ in $\lexab$. Let us show that $F\subset \homg lr,\om\la$. Let $x\in F\setminus M$. Then $wt(x,\lexab)=\om$. By Proposition \ref{p:1+1}, $\pi\chi(x,\lexab)=\om$. By (5)(b), $x\in \homg lr,\om\la$. Thus $F\subset \homg lr,\om\la$.
%Compact LOTS are sequentially compact, so a convergent sequence can be extracted from $M$.

Let us prove (8). Let $r\in \homg r,\om{\om_1}$. Take a strictly increasing sequence $\sq{x_\al}{\al<\om_1}\subset \homg l,\om{\om_1}$ converging to $r$. Let $F$ be the closure of $\sq{x_\al}{\al<\om_1}$ in $\lexab$. Then $F\subset \homg lr,\om{\om_1}$ and $F$ is homeomorphic to $\om_1+1$.
\end{proof}


\begin{proposition}\label{p:blex:2-1}
Let $f: X\to Y$ be a surjective increasing map between LOTS $X$ and $Y$.
If $f^{-1}(y)$ is a closed interval in $X$ for every $y\in Y$, then $f$ is continuous.
\end{proposition}
\begin{proof}
It suffices to show that $f^{-1}(\lL(y,-\infty))$ and $f^{-1}(\lR(y,+\infty))$ are open for every $y\in Y$.
Let $\ilcc[u,v]=f^{-1}(y)$. Then $f^{-1}(\lL(y,-\infty))=\lL(u,-\infty)$ and $f^{-1}(\lR(y,-\infty))=\lR(v,+\infty)$.
\end{proof}

For ordinals $\ka\leq \la$ denote by
\[
\pi_\ka^\la: \lexab \to \B^\ka,\ x\mapsto \restr x\ka.
\]
the projection.

The following statement follows from Proposition \ref{p:blex:2-1}.

\begin{proposition}\label{p:blex:2}
Let $\ka\leq \la$ be ordinals. The projection $\pi^\la_\ka: \lexab \to \B^\ka$ is continuous.
\end{proposition}


%\begin{proposition}\label{p:blex:t1}
\begin{theorem}\label{t:blex:t1}
Let $X$ be a zero-dimensional compact LOTS, $\la$ a regular cardinal.
The following conditions are equivalent:
\begin{enumerate}
\item
$X$ and $\lexab$ are order isomorphic;
\item
$X$ and $\lexab$ are homeomorphic;
%\item
%$\chi(x,X)=\la$ for all $x\in X$ and $X$ embeds in a product $\prod_{\al<\la}X_\al$ of spaces such that $\chi(X_\al)<\la$ for $\al<\la$;
\item
$\chi(x,X)=\la$ for all $x\in X$ and $X$ embeds in a product $\prod_{\al<\la}X_\al$ of spaces such that $\chi(X_\al)<\la$ for $\al<\la$;
\item
$\chi(x,X)=\la$ for all $x\in X$ and the set $M$ of jump points of the LOTS $X$ decomposes into a union $M=\bigcup_{\al<\la} M_\al$ of subsets such that $\chi(\cl{M_\al})<\la$ for $\al<\la$.
\end{enumerate}
\end{theorem}
\begin{proof}
(1) $\rarr$ (2). Obvious.

(2) $\rarr$ (3). By Proposition \ref{p:blex:2}, $\lexab$ embeds via the map $\diag_{\al<\la} \pi_\al^\la$ into the product $\prod_{\al<\la} \lex\al\B$. It remains to note that $\chi(\lex\al\B)\leq |\al|<\la$ for $\al<\la$.

(3) $\rarr$ (4). Assume $X\subset \prod_{\al<\la}X_\al$. We may assume that each $X_\al$ is compact. Denote $P_\be = \prod_{\al<\be}X_\al$ for $\be\leq\la$ and let $p_\be$ be the projection from $P_\la$ onto $P_\be$ for $\be<\la$.

Let $\be<\la$. Since $\la$ is a regular cardinal and $\chi(X_\al)<\la$ for $\la<\be$, we have $\ka'=\sup\{ \chi(X_\al) : \al < \be \}<\la$. Set $\ka=\ka' + |\be|$. Then $\chi(P_\be)\leq\ka<\la$.

For $(a,b)\in\jumps X$ set
\[
\te_{a,b}=\min \set{\al<\la: p_\al(\ilpc(-\infty,a]) \cap  p_\al(\ilcp[b,+\infty))=\es}.
\]
Set $M_\be=\set{a,b : (a,b)\in\jumps X,\ \te_{a,b}=\be}$ for $\be<\la$. Then $M=\bigcup_{\be<\la} M_\be$. Let us show that $\chi(\cl{M_\be})<\la$. Suppose otherwise. Then $\chi(x,\cl{M_\be})=\la$ for some $x\in \cl{M_\be}$. Let $L=\ilpp(-\infty,x)\cap M_\be$ and $R=\ilpp(x,+\infty)\cap M_\be$. Then $\chi(x,\cl{L})=\la$ or $\chi(x,\cl{R})=\la$. Assume, without loss of generality, that $\chi(x,\cl{R})=\la$. Let $\ka=\chi(p_\be(x), P_\be)$ and $F=p_\be^{-1}(p_\be(x))$. Since $\chi(x,\cl{R})=\la$, $F$ is a set of type $G_\ka$ and $\ka<\la$, there is $y>x$ with $\ilpp(x,y)\subset F$, and the set $\ilpp(x,y)\cap M_\be$ is infinite. There is $(a,b)\in\jumps X$ with $\sset{a,b}\subset \ilpp(x,y)\cap M_\be$. Then $\sset{a,b}\subset F$, and hence $p_\be(a)=p_\be(b)=x$, contradicting $p_\be(\ilpc(-\infty,a]) \cap  p_\be(\ilcp[b,+\infty))=\es$.

(4) $\rarr$ (1). Set $\cM_\al=\set{(a,b)\in \jumps X: \sset{a,b}\cap M_\al\neq\es}$ for $\al<\la$ and define the map
\[
\Om: \jumps X \to \la,\ (a,b) \mapsto \min\set{\al<\la: (a,b)\in \cM_\al}.
\]
Set $\te_{u,v}=\min  \Om(\jumps{\ilcc[u,v]})$ and $\cJ_{u,v}=\Om^{-1}(\te_{u,v})\cap \jumps{\ilcc[u,v]}$ for $u,v\in X$, $u<v$. Fix $(a_{u,v},b_{u,v})\in \cJ_{u,v}$ for every $u,v\in X$, $u<v$.

By induction on $\al<\la$, we construct increasing continuous maps $f_\al: X\to \B^\al$ such that the following conditions hold:
\begin{enumerate}
\item
$f_\al=\pi_\al^\be\circ f_\be$, where $\al<\be<\la$ and $\pi_\al^\be: \B^\be\to \B^\al$ is the projection;
%\item if $\al<\la$ is a limit ordinal, then 
\item
if $\al<\la$, 
$s\in \B^\al$ and $\ilcc[u,v]=f^{-1}_\al(s)$, then $f^{-1}_{\al+1}(s\fun 0)=\ilcc[u,a_{u,v}]$ and $f^{-1}_{\al+1}(s\fun 1)=\ilcc[b_{u,v},v]$.
\end{enumerate}
Base of the induction, $\al=0$. Then $\B^\al=\sset\es$; let $f_\al: X\to \sset\es$ be the identity map.

Inductive step, $0<\al<\la$. Consider the case when $\al$ is a limit ordinal. For $x\in X$ define $f_\al(x)\in \B^\al$: $f_\al(x)(\be)=f_{\be+1}(x)(\be)$ for $\be<\al$. Consider the case when $\al$ is a non-limit ordinal, $\al=\be+1$.
Define the map $g_\al: X\to \B$. Let $x\in X$. Since $\chi(x,X)=\la$ and $\chi(\B^\al)<\la$, we have $u<v$, where $\ilcc[u,v]=f^{-1}_\be(f_\be(x))$. Set
\[
g_\al(x) = \begin{cases}
0,&	x\in \ilcc[u,a_{u,v}],
\\
1,&	x\in \ilcc[b_{u,v},v].
\end{cases}
\]
Define $f_\al(x)\in \B^\al$:
\[
f_\al(x)(\ka) = \begin{cases}
f_\be(x)(\ka),&	\ka<\be,
\\
g_\al(x),&	\ka=\be.
\end{cases}
\]
The maps $f_\al$ for $\al<\la$ have been constructed.

Define the map $f: X\to \lexab$. For $x\in X$ set $f(x)(\al)=f_{\al+1}(x)(\al)$.
Since each $f_\al$ is a monotone map, $f$ is a monotone map.

For $s\in \lexab$ and $\al<\la$ denote 
\begin{align*}
\ilcc[u_\al^s,v_\al^s]&=f^{-1}_\al(\pi_\al^\la(s)), 
&
(a_\al^s,b_\al^s)&=(a_{u_\al^s,v_\al^s},b_{u_\al^s,v_\al^s}).
\end{align*}
Then
\begin{align*}
\ilcc[u_{\be+1}^s,v_{\be+1}^s] &=
\begin{cases}
\ilcc[u_{\be}^s,a_{\be}^s], &  s(\be)=0,
\\
\ilcc[b_{\be}^s,v_{\be}^s], &  s(\be)=1,
\end{cases}
\\
\ilcc[u_{\ga}^s,v_{\ga}^s] &= \bigcap_{\al<\ga}\ilcc[u_\al^s,v_\al^s],
\end{align*}
for $\be<\la$ and limit $\ga<\la$. Hence 
\[
f^{-1}(s) = \ilcc[u^s,v^s]=\bigcap_{\al<\la}\ilcc[u_\al^s,v_\al^s]\neq\es.
\]
Since $f^{-1}(s)\neq\es$ for all $s\in\lexab$, $f$ is surjective.
Since $f$ is a monotone surjection and $f^{-1}(s)$ is a closed interval for all $s\in\lexab$, Proposition \ref{p:blex:2-1} implies that $f$ is continuous.

Let us show that $f$ is a homeomorphism; for this it suffices to check that $f$ is injective. Suppose otherwise. Then $|f^{-1}(s)|>1$, and hence $u^s<v^s$ for some $s\in\lexab$.
%
Denote $\te^s=\te_{u^s,v^s}$ and $\te_\al=\te_{u_\al^s,v_\al^s}$ for $\al<\la$.
Then
\[
\ilcc[u^s,v^s] \subsetneq \ilcc[u^s_\be,v^s_\be] \subsetneq \ilcc[u^s_\al,v^s_\al]
\qquad\text{ and } \qquad \te_\al\leq \te_\be\leq  \te^s < \la
\]
for $\al< \be < \la$. 
Hence either $u^s_\al<u^s$ for all $\al<\la$, or $v^s<v^s_\al$ for all $\al<\la$. Assume, without loss of generality, that $u^s_\al<u^s$ for all $\al<\la$.
Since $\la$ is a regular cardinal and the sequence $\sq{\te_\al}{\al<\la}$ is bounded and increasing, the sequence $\sq{\te_\al}{\al<\la}$ stabilizes: for some $\te<\la$ and $\al_0<\la$, $\te_\al=\te$ for all $\al$ with $\al_0\leq \al < \la$. Set
\[
C=\set{ \al<\la: \al_0<\al,\ u^s_{\al_0} < u^s_{\al} < u^s_{\al+1} }
\]
Then $|C|=\la$, $u^s_{\al} < a^s_\al< b^s_\al=u^s_{\al+1}$, and $(a^s_\al,b^s_\al)\in \cM_\te$ for $\al\in C$.
Set $A=\set{a^s_\al:\al\in C}$, $B=\set{b^s_\al:\al\in C}$, and $S=M_\te\cap(A\cup B)$. Since $\sset{a^s_\al,b^s_\al}\cap M_\te\neq\es$ for $\al\in C$, we get $u^s\in \cl S$ and $\chi(u^s,\cl S)=\la$. Hence $u^s\in \cl{M_\te}$ and $\chi(u^s,\cl{M_\te})=\la$, contradicting $\chi(\cl{M_\te})<\la$.
\end{proof}

\begin{proposition}\label{p:blex:ct1}
Let $\la$ be a regular cardinal, $X$ a closed subset of $\lexab$.
The following conditions are equivalent:
\begin{enumerate}
\item
$X$ and $\lexab$ are order isomorphic;
\item
$\chi(x,X)=\la$ for all $x\in X$.
\end{enumerate}
\end{proposition}
\begin{proof}
(1) $\rarr$ (2) follows from Theorem \ref{t:blex:t1}.
Let us prove (2) $\rarr$ (1). Since, by Theorem \ref{t:blex:t1}, condition (3) of Theorem \ref{t:blex:t1} holds for $\lexab$ and $X\subset \lexab$, condition (3) of Theorem \ref{t:blex:t1} holds for $X$.
Hence $X$ and $\lexab$ are order isomorphic.
\end{proof}


\begin{proposition}\label{p:blex:3-1}
Let $X$ be a compact zero-dimensional LOTS. Then $w(X)=|\jumps X|$.
\end{proposition}
\begin{proof}
Since $X$ is a compact zero-dimensional LOTS, $\cU=\set{\ilpp(-\infty,a), \ilpp(b,+\infty): (a,b)\in \jumps X}$ forms a clopen subbase of the zero-dimensional compact space $X$. Hence $w(X)=|\jumps X|$.
\end{proof}


\begin{proposition}\label{p:blex:3}
Let $\la$ be a limit ordinal. Then $w(\lexab)=c(\lexab)=\sup\set{2^\al: \al<\la}$.
\end{proposition}
\begin{proof}
%Clearly, $|\jumps{\lexab}|$.
By Proposition \ref{p:blex:2}, $\lexab$ embeds via the map $\diag_{\al<\la} \pi_\al^\la$ into the product $\prod_{\al<\la} \lex\al\B$. Hence $w(\lexab)\leq \tau= \sup\set{2^\al: \al<\la}$.
For $\al<\la$, under the map $\pi^\la_\al$ the preimage of each point has nonempty interior. Hence $c(\lexab)\geq \tau$. It remains to note that $c(\lexab)\leq w(\lexab)$.
\end{proof}

\begin{theorem}\label{t:blex:1}
Let $\la$ be a regular cardinal, $X$ a compact zero-dimensional LOTS with $w(X)=\la$.
The following conditions are equivalent:
\begin{enumerate}
\item
$X$ and $\lexab$ are order isomorphic;
\item
$\chi(x,X)=\la$ for all $x\in X$.
\end{enumerate}
If either of these conditions holds, then $\la=\sup\set{2^\al: \al<\la}$.
\end{theorem}
\begin{proof}
By Proposition \ref{p:blex:3}, (1) implies $\la=\sup\set{2^\al: \al<\la}$.
(1) $\rarr$ (2) follows from Theorem \ref{t:blex:t1}.
Let us prove (2) $\rarr$ (1). Since $w(X)=\la$, $X$ embeds in $[0,1]^\la$, so condition (3) of Theorem \ref{t:blex:t1} holds for $X$. Hence $X$ and $\lexab$ are order isomorphic.
\end{proof}

The following statement follows from Theorem \ref{t:blex:1}.

\begin{cor}\label{c:blex:1}
{\rm ($\mathrm{CH}$)}
Let $X$ be a compact zero-dimensional LOTS.
The following conditions are equivalent:
\begin{enumerate}
\item
the LOTS $X$ is order isomorphic to $\lex {\om_1} \B$;
\item
$w(X)=\om_1$ and $\chi(x,X)=\om_1$ for all $x\in X$.
\end{enumerate}
\end{cor}

\begin{proposition}\label{p:blex:5}
%{\rm ($\mathrm{CH}$)}
$\lex {\om_1} \B$ does not embed homeomorphically in a homogeneous GO space.
\end{proposition}
\begin{proof}
Suppose $\lex {\om_1} \B$ embeds homeomorphically in a homogeneous GO space $X$, with $X\subset Y$ where $Y$ is a LOTS. Let $K\subset X$ be a compact subset homeomorphic to $\lex {\om_1} \B$.
By Theorem \ref{t:blex:t1}, $K$ is order isomorphic to $\lex {\om_1} \B$. Then $\cf_l(x,K)=\cf_r(x,K)=\om_1$ for some $x\in K$. Hence $\cf_l(x,X)=\cf_r(x,X)=\om_1$. By Proposition \ref{p:1+0}, $x$ is a WOB-point in $X$ and $\chi(x,X)=\om_1$. Hence $x$ is a $P$-point in $X$. By homogeneity of $X$, $X$ is a $P$-space, contradicting the fact that compact subsets of $P$-spaces are finite.
\end{proof}




\begin{theorem}\label{t:blex:2}
Let $\la$ be a regular cardinal, $a,b\in \lexab$, $a<b$.
The following conditions are equivalent:
\begin{enumerate}
\item
the LOTS $\ilcc[a,b]$ and $\lexab$ are order isomorphic;
\item
$\cf_r(a,\lexab)=\cf_l(b,\lexab)=\la$.
\end{enumerate}
\end{theorem}
\begin{proof}
The points $a$ and $b$ are not isolated in $\ilcc[a,b]$, so Proposition \ref{p:1+1.5} gives $\chi(a,\ilcc[a,b])=\cf_r(a,\lexab)$. For $x\in \ilpp(a,b)$, $\chi(x,\ilcc[a,b])=\chi(x,\lexab)=\la$.
Hence condition (2) is equivalent to $\chi(x,\ilcc[a,b])=\la$ for $x\in \ilcc[a,b]$.
Then (1) $\lrarr$ (2) follows from Proposition \ref{p:blex:ct1}.
\end{proof}

\begin{theorem}\label{t:blex:3}
Let $\la$ be a regular cardinal, $a,b,c,d\in \lexab$, $a<b$ and $c<d$.
The following conditions are equivalent:
\begin{enumerate}
\item
the LOTS $\ilcc[a,b]$ and $\ilcc[c,d]$ are order isomorphic and have no isolated points;
\item
$\cf_r(a,\lexab)=\cf_r(c,\lexab)\geq \om$ and
$\cf_l(b,\lexab)=\cf_l(d,\lexab)\geq \om$.
\end{enumerate}
\end{theorem}
\begin{proof}
Let 
$\al=\cf_r(a,\lexab)=\cf_r(c,\lexab)$ and
$\be = \cf_l(b,\lexab)=\cf_l(d,\lexab)$.

Clearly, (1) $\rarr$ (2). Let us prove (2) $\rarr$ (1).
Take $(a',b'),(c',d')\in \jumps{\lexab}$ such that $a< a'<b'< b$ and $c< c'<d'< d$.
By Proposition \ref{p:blex:1}, $\cf_l(a',\lexab)=\cf_r(b',\lexab)=\cf_l(c',\lexab)=\cf_r(d',\lexab)=\la$. Since $\ilcc[a,b]=\ilcc[a,a']\cup \ilcc[b',b]$ and $\ilcc[c,d]=\ilcc[c,c']\cup \ilcc[d',d]$, it suffices to prove that $\ilcc[a,a']$ is order isomorphic to $\ilcc[c,c']$ and $\ilcc[b',b]$ is order isomorphic to $\ilcc[d',d]$. This fact follows from the theorem being proved in the case $\al=\la$ or $\be=\la$. Let us prove (2) $\rarr$ (1) in the case $\al=\la$. The set $M=\set{x\in\lexab: \pi\chi(x,\lexab)=\la}$ is dense in $\lexab$. Take strictly increasing sequences $\sq{x_\te}{\te<\be}$ and $\sq{y_\te}{\te<\be}$ in $\lexab$ with $x_0=a$, $y_0=c$, $x_\te,y_\te\in M$ for $\te>0$, $\sq{x_\te}{\te<\be}$ cofinal in $\ilcp[a,b)$, and $\sq{y_\te}{\te<\be}$ cofinal in $\ilcp[c,d)$. Then $\cf_r(x_\te,\lexab)=\cf_l(x_{\te+1},\lexab)=\cf_r(y_\te,\lexab)=\cf_l(y_{\te+1},\lexab)=\la$ for $\te<\be$. By Theorem \ref{t:blex:2} there is an order isomorphism $f_\te: \ilcc[x_\te,x_{\te+1}]\to \ilcc[y_\te,y_{\te+1}]$ for $\te<\be$. Define an order isomorphism $f: \ilcc[a,b]\to \ilcc[c,d]$:
\[
f(x)=\begin{cases}
y_\te, & \text{if }x=x_\te\text{ for some }\te<\be,
\\
f_\te(x), & \text{if }x\in \ilpp(x_\te,x_{\te+1})\text{ for some }\te<\be,
\\
d, & \text{if }x=b
\end{cases}
\]
for $x\in \ilcc[a,b]$.
\end{proof}





%%% end file('ru//sec/blex.tex') %%%

\section{Homogeneous subspaces of the space of binary sequences} \label{sec:blexhome}
%%% begin file('ru//sec/blexhome.tex') %%%
\begin{proposition}\label{p:blexhome:1-1}
Let $\la$ be a regular cardinal and $x,y\in \lexab$.
The following conditions are equivalent:
\begin{enumerate}
\item
$f(x)=y$ for some $f\in\oip\lexab$;
\item
$\cf_r(x,\lexab)=\cf_r(y,\lexab)$ and
$\cf_l(x,\lexab)=\cf_l(y,\lexab)$.
\end{enumerate}
\end{proposition}
\begin{proof}
Let 
$\al=\cf_r(x,\lexab)=\cf_r(y,\lexab)$ and
$\be = \cf_l(x,\lexab)=\cf_l(y,\lexab)$.
By Proposition \ref{p:blex:1}, $\la\in\sset{\al,\be}$.
Clearly, (1) $\rarr$ (2). Let us prove (2) $\rarr$ (1).
If $0\in\sset{\al,\be}$, then $x=y$ and the statement is obvious.

Consider the case $1\in\sset{\al,\be}$. Consider the case $\al=1$. Then $x$ and $y$ are right jump points and $\be=\la$. Let $x'$ and $y'$ be the corresponding left jump points. Then $\cf_r(0_\la)=\cf_l(x')=\cf_r(x)=\cf_l(y')=\cf_r(y)=\cf_l(1_\la)=\la$. By Theorem \ref{t:blex:3} there exist order isomorphisms $f_1: \ilcc[0_\la,x']\to \ilcc[0_\la,y']$ and $f_2: \ilcc[x,1_\la]\to \ilcc[y,1_\la]$. For $u\in \lexab$, set
\[
f(u)=\begin{cases}
f_1(u), & u \leq x',
\\
f_2(u), & u \geq x.
\end{cases}
\]
Then $f\in\oip{\lexab}$ and $f(x)=y$.
Consider the case $\be=1$. Then $\cf_l(R_\la(x),\lexab)=\cf_l(R_\la(y),\lexab)=1$ and $\cf_r(R_\la(x),\lexab)=\cf_r(R_\la(y),\lexab)=\la$. Hence $f'(R_\la(x))=R_\la(y)$ for some $f'\in \oip{\lexab}$. We obtain $f(x)=y$ for $f=R_\la\circ f' \circ R_\la\in \oip\lexab$.

Consider the case when $\al$ and $\be$ are infinite cardinals.
Since $0_\la<x<1_\la$, $0_\la<y<1_\la$, $\al\geq \om$, and $\be\geq \om$, Theorem \ref{t:blex:3} gives order isomorphisms $f_1: \ilcc[0_\la,x]\to \ilcc[0_\la,y]$ and $f_2: \ilcc[x,1_\la]\to \ilcc[y,1_\la]$. For $u\in \lexab$, set
\[
f(u)=\begin{cases}
f_1(u), & u < x,
\\
y, & u=x,
\\
f_2(u), & u > x.
\end{cases}
\]
Then $f\in\oip{\lexab}$ and $f(x)=y$.
\end{proof}



\begin{proposition}\label{p:blexhome:1}
Let $\la$ be a regular cardinal and $\ka\leq\la$ a cardinal.
\begin{enumerate}
\item
If $x,y\in \homg l,\ka\la$ or $x,y\in \homg r,\ka\la$, then there is $f\in\oip X$ with $f(x)=y$.
\item
If $x\in \homg l,\ka\la$, $y\in \homg r,\ka\la$, or $x\in \homg r,\ka\la$, $y\in \homg l,\ka\la$, then there is $f\in\oim X$ with $f(x)=y$.
\item
If $x,y\in \homg lr,\ka\la$, then there is $f\in\oipm X$ with $f(x)=y$ and $f(\homg lr,\ka\la)=\homg lr,\ka\la$.
\item
if $f\in\oip X$, then $f(\homg l,\ka\la)=\homg l,\ka\la$, $f(\homg r,\ka\la)=\homg r,\ka\la$, and $f(\homg lr,\ka\la)=\homg lr,\ka\la$;
\item
$R_\la(\homg l,\ka\la)=\homg r,\ka\la$, $R_\la(\homg r,\ka\la)=\homg l,\ka\la$, and $R_\la(\homg lr,\ka\la)=\homg lr,\ka\la$;
\item
if $f\in\oim X$, then $f(\homg l,\ka\la)=\homg r,\ka\la$, $f(\homg r,\ka\la)=\homg l,\ka\la$, and $f(\homg lr,\ka\la)=\homg lr,\ka\la$.
\end{enumerate}
\end{proposition}
\begin{proof}
Items (4), (5) and (6) are obvious. Item (1) follows from Proposition \ref{p:blexhome:1-1}.

Let us prove (2). Let $x'=R_\la(x)$. By (1), there is $f'\in\oip \lexab$ with $f'(x')=y$. Let $f = f' \circ R_\la$. Then $f\in\oim \lexab$ and $f(x)=y$.

Let us prove (3). If $x,y\in \homg l,\ka\la$ or $x,y\in \homg r,\ka\la$, then (3) follows from (1) and (4).
If $x\in \homg l,\ka\la$, $y\in \homg r,\ka\la$, or $x\in \homg r,\ka\la$, $y\in \homg l,\ka\la$, then (3) follows from (2) and (6).
\end{proof}

The following proposition follows from Proposition \ref{p:blexhome:1}.

\begin{proposition}\label{p:blexhome:2}
Let $\la$ be a regular cardinal.
\begin{enumerate}
\item
The family
\begin{gather*}
\{ \homg l,0\la, \homg r,0\la, \homg l,1\la, \homg r,1\la, \homm\la \} \cup 
\\
\set{\homg l,\ka\la, \homg r,\ka\la: \ka<\la\text{ is a regular cardinal}}
\end{gather*}
is the family of orbits of the action of the group $\oip\lexab$ on $\lexab$.
\item
The family
\begin{gather*}
\{ \homg lr,0\la,  \homg lr,1\la, \homm\la \} \cup 
\\
\set{\homg lr,\ka\la: \ka<\la\text{ is a regular cardinal}}
\end{gather*}
is the family of orbits of the action of the group $\oipm\lexab$ on $\lexab$.
\end{enumerate}
\end{proposition}


\begin{theorem}\label{t:blexhome:1}
Let $\ka\leq\la$ be regular cardinals. The spaces $\homg l,\ka\la$, $\homg r,\ka\la$ are order-homogeneous LOTS. The space $\homg lr,\ka\la$ is a $\pm$-order-homogeneous LOTS.
\end{theorem}
\begin{proof}
By Proposition \ref{p:blex:1+1}(3)(b), the spaces $\homg l,\ka\la$, $\homg r,\ka\la$ and $\homg r,\ka\la$ are LOTS with the order inherited from $\lexab$. By Proposition \ref{p:blexhome:2}(1), the LOTS $\homg l,\ka\la$ and $\homg r,\ka\la$ are order-homogeneous. By Proposition \ref{p:blexhome:2}(2), the LOTS $\homg lr,\ka\la$ is $\pm$-order-homogeneous.
\end{proof}

\begin{cor}\label{c:blexhome:1}
Let $\ka\leq\la$ be regular cardinals. The spaces $\homg l,\ka\la$, $\homg r,\ka\la$ and $\homg lr,\ka\la$ are homogeneous LOTS.
\end{cor}






%%% end file('ru//sec/blexhome.tex') %%%

\section{Homogeneous extensions of LOTS} \label{sec:exos}
%%% begin file('ru//sec/exos.tex') %%%
For a LOTS $X$ denote 
\begin{align*}
\er(X) &= \set{x\in X: x\text{ is the minimum or a right jump point of }X},
\\
\el(X) &= \set{x\in X: x\text{ is the maximum or a left jump point of }X},
\\
\ea(X) &=  \el(X)\cup\er(X).
\end{align*}

Let $L$ and $R$ be LOTS such that $L$ has a minimum and no maximum, and $R$ has a maximum and no minimum.
Denote
\begin{align*}
\cD(X,L,R) &= (\el(X) \times R) \cup (\er(X) \times L),
\\
%\ex_1(X,Y) &= X \cup \left(\el(X) \times \nl \times Y\right) \cup \left(\er(X) \times \nr \times Y\right).
\ex(X,L,R) &= X \cup  \cD(X,L,R),
\\
\ex_0(X,L,R) &= X,
\\
\ex_n(X,L,R) &= \ex(\ex_{n-1}(X,L,R),L,R),
\\
\ex_\om(X,L,R) &= \bigcup_{n\in\om} \ex_n(X,L,R).
\end{align*}

We introduce a linear order on the set $\ex(X,L,R)$. On $X\subset \ex(X,L,R)$ the order is the same as on $X$. On $\cD(X,L,R)\subset X\times Y$ we use the lexicographic order. Let us compare $x\in X$ and $(z,y)\in \cD(X)$: $x<(z,y)$ if $x<z$, or $x=z$ and $(z,y)\in \el(X)\times R$; $x>(z,y)$ if $x>z$, or $x=z$ and $(z,y)\in \er(X)\times L$.


%For $(x,n)\in \cD(X,Y)$

The following statement follows from Proposition \ref{p:1+4} and the construction of the spaces $\ex(X,L,R)$ and $\ex_\om(X,L,R)$.

\begin{proposition}\label{p:exos:1}
Let $X$ be a compact LOTS, $L$ and $R$ locally compact LOTS, $L$ having a minimum and no maximum, and $R$ having a maximum and no minimum. Then $X$ embeds in $\ex(X,L,R)$, $\ex(X,L,R)$ is compact, and $\ex_\om(X,L,R)$ is $\sigma$-compact.
\end{proposition}

Below we shall omit the parameters $L$ and $R$: we shall write $\cD(X)$ instead of $\cD(X,L,R)$, and similarly for the other notation.


For $x\in \ea(X)$ and $M\subset \ea(X)$ set
\begin{align*}
T_{x}(X) &= \begin{cases}
 R,& x\in \el(X),
\\
 L,& x\in \er(X),
\end{cases}
&
\cD_{x}(X) &= \sset x \times T_x(X),
\\
\cD_{M}(X) &= \bigcup_{x\in M} \cD_{x}(X),
&
\ex(X;M) &= X \cup \cD_{M}(X).
\end{align*}


\begin{statement}\label{s:exos:1}
Let $M\subset \ea(X)$.
Then $\ex_\om(\ex(X;M))=\ex_\om(X)$.
\end{statement}
\begin{proof}
By induction on $n\in\om$ we define sets $M_n\subset \ea(X_n)$, where $X_n=\ex_n(X)$.
Set $M_0=M$. Set $M_{n+1}=\ea(\cD_{M_n}(X_n))$ for $n>0$. Set $Z_n=\ex(X_n;M_n)$ for $\nom$. Then $Z_0=\ex(X;M)$, $Z_{n+1}=\ex(Z_n)$, and $X_n\subset Z_n\subset X_{n+1}$. Hence $\ex_\om(X)=\bigcup_\nom X_n=\bigcup_\nom Z_n=\ex_\om(Z_0)$.
\end{proof}


\begin{statement}\label{s:exos:2}
If $R$ is order reversing isomorphic to $L$, then every $f\in\oipm X$ extends to some $\hat f\in \oipm{\ex_\om(X)}$.
\end{statement}
\begin{proof}
Let $\ph: L\to R$ and $\psi: R\to L$ be order reversing isomorphisms such that $\ph\circ \psi$ is the identity map of $L$ onto itself.
\begin{claim}\label{cl:exos:1-1}
Let $Y$ be a LOTS. Every $f\in\oipm Y$ extends to some $\hat f\in \oipm{\ex(Y)}$.
\end{claim}
\begin{proof}
Suppose $f\in \oip Y$. Set
\[
\hat f(x) = \begin{cases}
f(x), & x\in Y,
\\
(f(y),z), & x=(y,z)\in \cD(Y).
\end{cases}
\]
Then $\hat f\in \oip{\ex(Y)}$.

Suppose $f\in \oim Y$. Set
\[
\hat f(x) = \begin{cases}
f(x), & x\in Y,
\\
(f(y),\ph(z)), & x=(y,z)\in \er(Y) \times L,
\\
(f(y),\psi(z)), & x=(y,z)\in \el(Y) \times R.
\end{cases}
\]
Then $\hat f\in \oim{\ex(Y)}$.
\end{proof}
By Fact \ref{cl:exos:1-1}, there exist $f_n\in\oipm{\ex_n(X)}$ for $\nom$ such that $f_0=f$ and $f_n=\restr{f_{n+1}}{\ex_n(X)}$ for $\nom$. Define $\hat f: \ex_\om(X)\to \ex_\om(X)$ by $\hat f(x)=f(x)$ for $x\in X$, and $\hat f(x)=f_n(x)$ for $x\in \ex_n(X)\setminus \ex_{n-1}(X)$.
\end{proof}


\begin{theorem}\label{t:exos:1}
Let $X$, $L$ and $R$ be LOTS.
Let $L$ have a minimum and no maximum, and let $L$ and $\rev R$ be order isomorphic.
Suppose the following conditions hold:
\begin{enumerate}
\item
if $x,y\in X\setminus \ea(X)$, then $f(x)=y$ for some $f\in\oipm X$;
\item
the LOTS $\ilpc(-\infty,a] + R + \ilpp(a,+\infty)$ and $\ilpp(-\infty,b) + L + \ilcp[b,+\infty)$ are order isomorphic to $X$ for $a\in \el(X)$ and $b\in \er(X)$.
\end{enumerate}
Then $\ex_\om(X,L,R)$ is a $\pm$-order-homogeneous LOTS.
\end{theorem}
\begin{proof}
Condition (2) is equivalent to $\ex(X;\sset x)$ being order isomorphic to $X$ for $x\in\ea(X)$.
\begin{claim}\label{cl:exos:1}
Let $u\in \ex_\om(X)$. There is $Y\subset \ex_\om(X)$ such that $X\subset Y$, $\ex_\om(Y)=\ex_\om(X)$, $Y$ is order isomorphic to $X$, and $u\in Y\setminus \ea(Y)$.
\end{claim}
\begin{proof}
For some $n\in\om$,
\begin{align*}
u &= (x,u_1,...,u_n)\in \ex_n(X).
\end{align*}
Set $x_0=x$, $x_1=(x,u_1)$, ..., $x_n=(x,u_1,...,u_n)=u$. Define $X_i\subset \ex_\om(X)$ for $i=0,1,...,n$ such that $x_i\in X_i$ for $i\leq n$ and $x_i \in \ea(X_i)$ for $i<n$. Set $X_0=X$ and $X_i=\ex(X_{i-1};\sset{x_{i-1}})$ for $i>0$. By Statement \ref{s:exos:1}, $\ex_\om(X_i) = \ex_\om(X)$, and by condition (2), $X_i$ is order isomorphic to $X$ for $i=0,1,...,n$. If $u=x_n\notin \ea(X_n)$, set $Y=X_n$. If $u\in \ea(X_n)$, set $Y=\ex(X_{n};\sset{u})$. Then $u\notin \ea(Y)$. By Statement \ref{s:exos:1} and (2), $Y$ is order isomorphic to $X$ and $\ex_\om(Y)=\ex_\om(X)$.
\end{proof}
Let $x,y\in\ex_\om(X)$. By Fact \ref{cl:exos:1}, there is $Y'\subset \ex_\om(X)$ such that $X\subset Y'$, $x\in Y'\setminus \ea(Y')$, $\ex_\om(Y')=\ex_\om(X)$, and $Y'$ is order isomorphic to $X$. By Fact \ref{cl:exos:1}, there is $Y\subset \ex_\om(Y')$ such that $Y'\subset Y$, $y\in Y\setminus \ea(Y)$, $\ex_\om(Y)=\ex_\om(Y')=\ex_\om(X)$, and $Y$ is order isomorphic to $Y'$ and to $X$.

Thus $x,y\in Y\setminus \ea(Y)$, $\ex_\om(Y)=\ex_\om(X)$, and $Y$ is order isomorphic to $X$.
By condition (1), $f(x)=y$ for some $f\in\oipm Y$. By Statement \ref{cl:exos:1-1}, $f$ extends to some $\hat f\in \oipm{\ex_\om(Y)}$. Then $\hat f(x)=y$ for $\hat f\in \oipm{\ex_\om(X)}$. Hence the LOTS $\ex_\om(X)$ is $\pm$-order-homogeneous.
\end{proof}






%%% end file('ru//sec/exos.tex') %%%

\section{Large compacta in homogeneous LOTS} \label{sec:bcs}
%%% begin file('ru//sec/bcs.tex') %%%

Set $\sS=(\lex{\om_1}\B\setminus \homm{\om_1})\cup (\homm{\om_1}\times\B)$. We define a linear order on $\sS$. The order on $\lex{\om_1}\B\setminus \homm{\om_1}$ is inherited from $\lex{\om_1}\B$, and the order on $\homm{\om_1}\times\B$ is lexicographic. Let $x\in \lex{\om_1}\B$ and $(y,a)\in \homm{\om_1}\times\B$. Then $x<(y,a)$ if and only if $x<y$, and $x>(y,a)$ if and only if $x>y$.

Define the map $\pi: \sS\to \lex{\om_1}\B$ by $\pi(x)=x$ for $x\in \lex{\om_1}\B\setminus \homm{\om_1}$, and $\pi((y,a))=y$ for $(y,a)\in \homm{\om_1}\times\B$. The map $\pi$ is a continuous increasing map of compact LOTS.

Define the order reversing automorphism $\Rs: \sS\to\sS$ by $\Rs(x)=R_{\om_1}(x)$ for $x\in \lex{\om_1}\B\setminus \homm{\om_1}$, $\Rs((y,0))=(R_{\om_1}(y),1)$, and $\Rs((y,1))=(R_{\om_1}(y),0)$ for $y\in  \homm{\om_1}$. Note that $\pi\circ R_{\om_1}=\Rs\circ \pi$.

For cardinals $\nu$ and $\mu$ set
\begin{align*}
\sS_{\nu,\mu} &= \set{x\in\sS: \cf_l(x,\sS)=\nu,\ \cf_r(x,\sS)=\mu},
&
\wS_{\nu,\mu} &= \sS_{\nu,\mu} \cup \sS_{\mu,\nu}.
\end{align*}

Note that $\mns$ is the minimum of $\sS$, and $\mxs$ is the maximum of $\sS$. From the construction of the space $\sS$ and Propositions \ref{p:blex:1+1} and \ref{p:1+1} we get
\begin{align*}
\sS_{0,\om_1}&=\sset{\mns},
\qquad
\sS_{\om_1,0}=\sset{\mxs},
\qquad
\wS_{0,\om_1}=\sset{\mns,\mxs},
\\
\sS_{1,\om_1}&=\set{x\in\sS : x\text{ is a right jump point}},
\\
\sS_{\om_1,1}&=\set{x\in\sS : x\text{ is a left jump point}},
\\
\wS_{1,\om_1}&=\set{x\in\sS : x\text{ is a jump point}},
\\
\wS_{\om,\om_1}&=\set{x\in\sS : \pi\chi(x,\lexab)=\om}
\end{align*}
and that the families
\begin{align*}
\cS &=
\{\,
\sS_{0,\om_1}, \sS_{\om_1,0}, \sS_{\om,\om_1}, \sS_{\om_1,\om}, \sS_{1,\om_1}, \sS_{\om_1,1}
\,\},
\\
\wcS &=
\{\,
\wS_{0,\om_1}, \wS_{\om,\om_1}, \wS_{1,\om_1}
\,\},
\end{align*}
are partitions of the set $\sS$.
%
%
Then the following statement follows from Propositions \ref{p:1+1} and \ref{p:blex:3-1}.

\begin{proposition}\label{p:bsc:1-1}
$\chi(x,\sS)=\om_1$ for all $x\in \sS$, and $w(\sS)=|\sS|=2^{\om_1}$.
\end{proposition}


\begin{proposition}\label{p:bsc:1}
Let $a,b,c,d\in \sS$, $a<b$ and $c<d$.
The following conditions are equivalent:
\begin{enumerate}
\item
the LOTS $\ilcc[a,b]$ and $\ilcc[c,d]$ are order isomorphic and have no isolated points;
\item
$\cf_r(a,\sS)=\cf_r(c,\sS)\geq \om$ and
$\cf_l(b,\sS)=\cf_l(d,\sS)\geq \om$.
\end{enumerate}
\end{proposition}
\begin{proof}
Let
$\al=\cf_r(a,\sS)=\cf_r(c,\sS)$ and
$\be = \cf_l(b,\sS)=\cf_l(d,\sS)$.
Clearly, (1) $\rarr$ (2). Let us prove (2) $\rarr$ (1). Since $\cf_r(\pi(a),\lexab)=\cf_r(\pi(c),\lexab)=\al$ and
$\cf_l(\pi(b),\lexab)=\cf_l(\pi(d),\lexab)=\be$, Theorem \ref{t:blex:3} gives an order isomorphism $f': \ilcc[\pi(a),\pi(b)]\to \ilcc[\pi(c),\pi(d)]$. Define an order isomorphism $f: \ilcc[a,b]\to\ilcc[c,d]$ by setting
\[
f(u)=\begin{cases}
f'(u), & u\in \homm{\om_1},
\\
(f'(x),y), & u\notin \homm{\om_1}\text{ and }u=(x,y),
\end{cases}
\]
for $u\in \ilcc[a,b]$.
\end{proof}

\begin{proposition}\label{p:bsc:1+1}
If $(a,b)\in \jumps\sS$, then $\ilcc[\mns,a]$ and $\ilcc[b,\mxs]$ are order isomorphic to $\sS$.
\end{proposition}
\begin{proof}
Since $\cf_r(\mns,\cS)=\cf_l(a,\cS)=\cf_r(b,\cS)=\cf_l(\mxs,\cS)=\om_1$, Proposition \ref{p:bsc:1} implies that the LOTS $\ilcc[\mns,a]$, $\ilcc[b,\mxs]$ and $\sS$ are order isomorphic.
\end{proof}


\begin{proposition}\label{p:bsc:2}
Let $x,y\in \sS$.
The following conditions are equivalent:
\begin{enumerate}
\item
$f(x)=y$ for some $f\in\oip\sS$;
\item
$\cf_r(x,\sS)=\cf_r(y,\sS)$ and
$\cf_l(x,\sS)=\cf_l(y,\sS)$.
\end{enumerate}
\end{proposition}
\begin{proof}
Let
$\al=\cf_r(x,\sS)=\cf_r(y,\sS)$ and
$\be = \cf_l(x,\sS)=\cf_l(y,\sS)$.
Clearly, (1) $\rarr$ (2). Let us prove (2) $\rarr$ (1).
If $0\in\sset{\al,\be}$, then $x=y$ and the statement is obvious.

Consider the case $1\in\sset{\al,\be}$. Consider the case $\al=1$. Then $x$ and $y$ are right jump points and $\be=\om_1$. Let $x'$ and $y'$ be the corresponding left jump points. By Proposition \ref{p:bsc:1+1}, there exist order isomorphisms $f_1: \ilcc[\mns,x']\to \ilcc[\mns,y']$ and $f_2: \ilcc[x,\mxs]\to \ilcc[y,\mxs]$. For $u\in \sS$, set
\[
f(u)=\begin{cases}
f_1(u), & u \leq x',
\\
f_2(u), & u \geq x.
\end{cases}
\]
Then $f\in\oip{\sS}$ and $f(x)=y$.
Consider the case $\be=1$. Then $\cf_l(\Rs(x),\sS)=\cf_l(\Rs(y),\sS)=1$ and $\cf_r(\Rs(x),\sS)=\cf_r(\Rs(y),\sS)=\la$. Hence $f'(\Rs(x))=\Rs(y)$ for some $f'\in \oip{\sS}$. We obtain $f(x)=y$ for $f=\Rs\circ f' \circ \Rs\in \oip\sS$.

Consider the case when $\al$ and $\be$ are infinite cardinals.
Since $\mns<x<\mxs$, $\mns<y<\mxs$, $\al\geq \om$, and $\be\geq \om$, Proposition \ref{p:bsc:1} gives order isomorphisms $f_1: \ilcc[\mns,x]\to \ilcc[\mns,y]$ and $f_2: \ilcc[x,\mxs]\to \ilcc[y,\mxs]$. For $u\in \sS$, set
\[
f(u)=\begin{cases}
f_1(u), & u < x,
\\
y, & u=x,
\\
f_2(u), & u > x.
\end{cases}
\]
Then $f\in\oip{\sS}$ and $f(x)=y$.
\end{proof}

\begin{proposition}\label{p:bsc:3}
\begin{enumerate}
\item
The family $\cS$
is the family of orbits of the action of the group $\oip\sS$ on $\sS$.
\item
The family $\wS$ is the family of orbits of the action of the group $\oipm\sS$ on $\sS$.
\end{enumerate}
\end{proposition}
\begin{proof}
Since $\cS$ is a partition of the space $\sS$, item (1) follows from Proposition \ref{p:bsc:2}.
Item (2) follows from the fact that $\Rs\in\oim\sS$, $\wS=\set{M\cup \Rs(M): M\in\sS}$, and item (1).
\end{proof}



\begin{proposition}\label{p:bsc:4}
The following LOTS are order isomorphic to $\sS$: $\sS + \sS$, $\om \times \sS + \sS$, $\sS+ \rev (\om \times \sS)$.
\end{proposition}
\begin{proof}
Let $(a,b)\in \jumps \sS$. By Proposition \ref{p:bsc:1+1}, the LOTS $\ilcc[\mns,a]$, $\ilcc[b,\mxs]$ and $\sS$ are order isomorphic. Hence $\sS$ and $\sS+\sS$ are order isomorphic.

Take a sequence of jumps $(a_i,b_i)_{i\in\om}\subset \jumps \sS$ such that $b_i<a_{i+1}$ for $i\in\om$. Let $b_{-1}=\mns$ and $b=\sup \set{b_i:i\in\om}$. Then $\cf_r(b_{i-1},\sS)=\cf_l(a_i,\sS)=\om_1$ for $i\in\om$, and $\cf_l(b,\sS)=\cf_r(\mxs,\sS)=\om_1$. By Proposition \ref{p:bsc:1}, the LOTS $\ilcc[b_{i-1},a_i]$ for $i\in\om$ and $\ilcc[b,\mxs]$ are order isomorphic to $\sS$. Hence $\ilcp[\mns,b)$ is order isomorphic to $\om\times \sS$. Since $\sS=\ilcp[\mns,b)+\ilcc[b,\mxs]$, $\sS$ is order isomorphic to $\om\times \sS + \sS$.

Since $\sS$ is order isomorphic to $\rev \sS$, $\sS+ \rev(\om\times \sS)$ is order isomorphic to $\rev(\om\times \sS + \sS)$. Hence $\sS+ \rev(\om\times \sS)$ is order isomorphic to $\sS$.
\end{proof}




\begin{theorem}\label{t:bsc:1}
The LOTS $\ex_\om(\sS, \om\times \sS, \rev(\om\times \sS))$ is $\pm$-order-homogeneous.
\end{theorem}
\begin{proof}
Let us check that the conditions of Theorem \ref{t:exos:1} hold. Since $\sS\setminus \ea(\sS)=\wS_{\om,\om_1}$, condition (1) of Theorem \ref{t:exos:1} follows from Proposition \ref{p:bsc:3}(2).


Let us check condition (2). Let $L=\om\times \sS$ and $R=\rev(\om\times \sS)$, $a\in \el(X)$ and $b\in \er(X)$.

If $a=\max \sS$, then $A=\ilpc(-\infty,a] + R + \ilpp(a,+\infty)=\sS + \rev(\om\times \sS)$. By Proposition \ref{p:bsc:4}, $A$ is order isomorphic to $\sS$. If $a$ is a left jump point, then $A= \ilcc[\mns,a] + \rev(\om\times \sS) + \ilcc[a',\mxs]$, where $a'=\min \ilpp(a,+\infty)$. By Proposition \ref{p:bsc:1+1}, $A$ is order isomorphic to $\sS + \rev(\om\times \sS) + \sS$. By Proposition \ref{p:bsc:4}, $\sS + \rev(\om\times \sS)$ is order isomorphic to $\sS$. Hence $A$ is order isomorphic to $\sS + \sS$. By Proposition \ref{p:bsc:4}, $A$ is order isomorphic to $\sS$.

If $b=\min \sS$, then $B=\ilpp(-\infty,b) + L + \ilcp[b,+\infty)=\om\times \sS + \sS$. By Proposition \ref{p:bsc:4}, $B$ is order isomorphic to $\sS$. If $b$ is a right jump point, then $B= \ilcc[\mns,b'] + \om\times \sS + \ilcc[b,\mxs]$, where $b'=\max \ilpp(-\infty,b)$. By Proposition \ref{p:bsc:1+1}, $A$ is order isomorphic to $\sS + \om\times \sS + \sS$. By Proposition \ref{p:bsc:4}, $\om\times \sS+\sS$ is order isomorphic to $\sS$. Hence $B$ is order isomorphic to $\sS + \sS$. By Proposition \ref{p:bsc:4}, $B$ is order isomorphic to $\sS$.

By Theorem \ref{t:exos:1}, $\ex_\om(\sS, \om\times \sS, \rev(\om\times \sS))$ is $\pm$-order-homogeneous.
\end{proof}
























%%% end file('ru//sec/bcs.tex') %%%

\section{Examples and questions} \label{sec:qe}
%%% begin file('ru//sec/qe.tex') %%%

A semitopological group $G$ that is a LOTS and not a $P$-space is metrizable (Theorem \ref{t:main:3}). In this statement the condition that $G$ is not a $P$-space cannot be dropped, and the condition that $G$ is a LOTS cannot be weakened to $G$ being a GO space.

The following statement follows from Proposition \ref{p:h:tg1}.
\begin{example}\label{e:1}
The group $\lex{\om_1}\R$ with the order topology is a non-metrizable LOTS group that is a $P$-space.
\end{example}

\begin{example}\label{e:2}
The Sorgenfrey line is a non-metrizable, first-countable GO paratopological group.
\end{example}

A semitopological group $G$ that is a GO space and not a $P$-space is submetrizable, paracompact, and has countable character (Theorem \ref{t:main:2}). In this statement the condition that $G$ is a semitopological group cannot be replaced by the condition that $G$ is a homogeneous space.

For a regular ordinal $\la$ denote $X_\la=\homg lr,\om\la$.

\begin{example}\label{e:3}
Let $\la$ be a regular ordinal.
The space $X_\la$ is a homogeneous countably compact LOTS with $\chi(X_\la)=\la$ and $\pi\chi(X_\la)=\om$.
\end{example}
\begin{proof}
By Corollary \ref{c:blexhome:1}, the space $X_\la$ is homogeneous. By Proposition \ref{p:blex:1+1}, $X_\la$ is a countably compact LOTS with $\pi\chi(X_\la)=\om$ and $\chi(X_\la)=\la$.
\end{proof}

By Theorem \ref{t:main:1}, if $\la>\om_1$ is a regular cardinal, then compact subsets of $X_\la$ are first countable.

\begin{example}\label{e:4}
The space $X_{\om_1}$ is a homogeneous countably compact LOTS with $\chi(X_{\om_1})=\om_1$, $\pi\chi(X_{\om_1})=\om$, and the transfinite line $\om_{1}+1$ embeds in $X_{\om_1}$.
\end{example}
\begin{proof}
By Proposition \ref{p:blex:1+1}, the transfinite line $\om_{1}+1$ embeds in $X_{\om_1}$.
\end{proof}

The LOTS $\lex {\om_1} \B$ does not embed in a homogeneous GO space (Proposition \ref{p:blex:5}).
Let $\Hs=\ex_\om(\sS, \om\times \sS, \rev(\om\times \sS))$.

\begin{example}\label{e:5}
The space $\Hs$ is a homogeneous $\sigma$-compact LOTS with $\chi(\Hs)=\om_1$ and $\pi\chi(\Hs)=\om$, and the compact space $\sS$ embeds in $\Hs$, with $w(\sS)=c(\sS)=|\sS|=2^{\om_1}$ and $\chi(x,\sS)=\om_1$ for all $x\in\sS$.
\end{example}
\begin{proof}
By Theorem \ref{t:bsc:1}, the space $\Hs$ is homogeneous. By Proposition \ref{p:exos:1}, $\Hs$ is $\sigma$-compact and contains the compact set $\sS$. By Proposition \ref{p:bsc:1-1}, $w(\sS)=c(\sS)=|\sS|=2^{\om_1}$ and $\chi(x,\sS)=\om_1$ for all $x\in\sS$.
By Theorem \ref{t:main:1}, $\chi(\Hs)=\om_1$ and $\pi\chi(\Hs)=\om$.
\end{proof}

\begin{question}\label{q:1-1}
Can the compact space $\sS$ be embedded in a homogeneous countably compact LOTS? Let $K$ be a compact subset of some homogeneous countably compact LOTS. Is it true that $\chi(x,K)\leq\om$ for some $x\in K$?
\end{question}


\begin{example}\label{e:6} \cite{Banakh2008}
The compact LOTS `double arrow' admits the structure of a right-topological group.
\end{example}

\begin{question}\label{q:1}
Does the LOTS $X_{\om_1}$ admit the structure of a right-topological group?
\end{question}

\begin{question}\label{q:1+1}
Does the LOTS $\Hs$ admit the structure of a right-topological group?
\end{question}


\begin{question}\label{q:2}
Let $X$ be a power homogeneous GO (LOTS) space.
Is it true that every compact subset of $X$ has tightness at most $\om_1$?
\end{question}

\begin{question}\label{q:3}
Let $X$ be a monotonically normal ($p$-)homogeneous space.
Is it true that every compact subset of $X$ has tightness at most $\om_1$?
\end{question}



There are many examples of separable, non-first-countable, non-metrizable monotonically normal topological groups \cite{GartsideReznichenko2000nmfs}. Such groups are stratifiable spaces
\cite[Theorem 24]{GartsideReznichenko2000nmfs}.

\begin{question}\label{q:4}
Let $G$ be a separable monotonically normal paratopological (semitopological, quasitopological) group. Is it true that $G$ is a stratifiable space?
\end{question}

\begin{question}[{\cite[Problem 6]{Collins1996}}]\label{q:5}
Let $G$ be a monotonically normal topological group with countable pseudocharacter. Is it true that $G$ is a stratifiable space?
\end{question}


\begin{question}\label{q:6}
Let $G$ be a monotonically normal paratopological (semitopological, quasitopological) group with countable pseudocharacter. Is it true that $G$ is a stratifiable space?
\end{question}

Monotonically normal topological vector spaces are stratifiable spaces \cite[Theorem 1]{Shkarin2004}.

\begin{question}[{\cite[Remarks (2)]{Shkarin2004}}]\label{q:7}
Let $G$ be a monotonically normal topological group and let $G$ contain a convergent sequence. Is it true that $G$ is a stratifiable space?
\end{question}

\begin{question}\label{q:8}
Let $G$ be a monotonically normal paratopological (semitopological, quasitopological) group and let $G$ contain a convergent sequence. Is it true that $G$ is a stratifiable space?
\end{question}

Monotonically normal paratopological groups are hereditarily paracompact \cite[Theorem 2.6]{BuzyakovaVural2014}, and hence compact subsets of such groups have countable tightness \cite[Theorem 3.5]{WilliamsZhou1998} (Theorem \ref{t:main:3+1}). Compact subsets of monotonically normal topological groups are metrizable \cite[Theorem 2a]{Gartside1998}.

\begin{question}\label{q:9}
Let $G$ be a monotonically normal paratopological (semitopological, quasitopological) group. Is it true that compact subsets of $G$ are metrizable (separable, first countable)?
\end{question}





%%% end file('ru//sec/qe.tex') %%%

%\input sec/text.tex


\begin{thebibliography}{10}
\expandafter\ifx\csname url\endcsname\relax
  \def\url#1{\texttt{#1}}\fi
\expandafter\ifx\csname urlprefix\endcsname\relax\def\urlprefix{URL }\fi
\expandafter\ifx\csname href\endcsname\relax
  \def\href#1#2{#2} \def\path#1{#1}\fi

\bibitem{Arhangelskii_vanMill}
A.~Arhangel’skii, J.~v. Mill, Topological homogeneity, in: Recent Progress in
  General Topology III, Springer, 2014, pp. 1--68.

\bibitem{arh67}
A.~Arhangelskii, Groupes topologiques extremalement discontinus, C.R.Acad. Sci.
  Paris 265 (1967) 822--825.

\bibitem{Arhangelskii2018}
A.~Arhangel'skii, Remainders of extremally disconnected spaces and related
  objects, Topology and its Applications 249 (2018) 9--18.
\newblock \href {https://doi.org/10.1016/j.topol.2018.09.001}
  {\path{doi:10.1016/j.topol.2018.09.001}}.

\bibitem{rezn2020}
E.~Reznichenko,
  \href{https://www.sciencedirect.com/science/article/pii/S016686412030345X}{Homogeneous
  subspaces of products of extremally disconnected spaces}, Topology and its
  Applications 284 (2020) 107403.
\newblock \href {https://doi.org/10.1016/j.topol.2020.107403}
  {\path{doi:10.1016/j.topol.2020.107403}}.
\newline\urlprefix\url{https://www.sciencedirect.com/science/article/pii/S016686412030345X}

\bibitem{Groznova2023ru}
А.~Ю. Грознова, Об однородности
  произведений топологических пространств,
  Математические заметки 113~(2) (2023) 171--181
  [A.~Yu.~Groznova, On the homogeneity of products of topological spaces, Mathematical Notes 113(2) (2023) 171--181 (in Russian)].
\newblock \href {https://doi.org/10.4213/mzm13634}
  {\path{doi:10.4213/mzm13634}}.

\bibitem{Arhangelskii2005}
A.~Arhangel’skii, Homogeneity of powers of spaces and the character,
  Proceedings of the American Mathematical Society 133~(7) (2005) 2165--2172.

\bibitem{BuzyakovaVural2014}
R.~Z. Buzyakova, C.~Vural, Stationary sets in topological and paratopological
  groups, Houston journal of mathematics 40~(1) (2014) 267--273.

\bibitem{GartsideReznichenko2000nmfs}
P.~Gartside, E.~Reznichenko, Near metric properties of function spaces,
  Fundamenta Mathematicae 164~(2) (2000) 97--114.

\bibitem{WilliamsZhou1998}
S.~W. Williams, H.~Zhou, Order-like structure of monotonically normal spaces,
  Commentationes Mathematicae Universitatis Carolinae 39~(1) (1998) 207--217.

\bibitem{Gartside1998}
P.~Gartside, Nonstratifiability of topological vector spaces, Topology and its
  Applications 86~(2) (1998) 133--140.

\bibitem{BennettBurkeLutzer2013}
H.~Bennett, D.~Burke, D.~Lutzer, Choban operators in generalized ordered
  spaces, Topology and its Applications 160~(4) (2013) 633--647.

\bibitem{gru1984}
G.~Gruenhage, Generalized metric spaces, in: Handbook of set-theoretic
  topology, Elsevier Science Publishers B. V., 1984, pp. 423--502.
\newblock \href {https://doi.org/10.1016/B978-0-444-86580-9.50013-6}
  {\path{doi:10.1016/B978-0-444-86580-9.50013-6}}.

\bibitem{JuhaszNyikosSzentmiklossy2005}
I.~Juh{\'a}sz, P.~Nyikos, Z.~Szentmikl{\'o}ssy, Cardinal restrictions on some
  homogeneous compacta, Proceedings of the American Mathematical Society
  133~(9) (2005) 2741--2750.

\bibitem{Hodel1984handbook}
R.~Hodel, {Cardinal functions I}, in: Handbook of set-theoretic topology,
  Elsevier, 1984, pp. 1--61.

\bibitem{Juhasz1992}
I.~Juh{\'a}sz, Cardinal functions, in: Recent Progress in General Topology,
  Amsterdam: Elsevier Science Publishers, 1992, pp. 417--441.

\bibitem{Nagata1985}
J.~Nagata, Modern general topology, Elsevier, 1985.

\bibitem{EngelkingBookRu}
Р. Энгелькинг, Общая топология, М.: Мир, 1986
[R.~Engelking, General Topology, Mir, Moscow, 1986 (Russian translation); originally published as: R.~Engelking, General Topology, PWN, Warsaw, 1977].

\bibitem{Arhangelskii2004}
A.~Arhangel’skii, A weak algebraic structure on topological spaces and
  cardinal invariants, Topology Proc 28~(1) (2004) 1--18.

\bibitem{Reznichenko2023-1}
E.~Reznichenko, {Metrizability of CHART groups}, Topology and its Applications
  326 (2023) 108408.
\newblock \href {https://doi.org/10.1016/j.topol.2023.108408}
  {\path{doi:10.1016/j.topol.2023.108408}}.

\bibitem{arh-rezn2005}
A.~Arhangel'skii, E.~Reznichenko, Paratopological and semitopological groups
  versus topological groups, Topology and its Applications 151~(1-3) (2005)
  107--119.

\bibitem{at2009}
A.~Arhangel'skii, M.~Tkachenko,
  \href{https://doi.org/10.2991/978-94-91216-35-0}{{Topological Groups and
  Related Structures}}, Atlantis Press, 2008.
\newblock \href {https://doi.org/10.2991/978-94-91216-35-0}
  {\path{doi:10.2991/978-94-91216-35-0}}.
\newline\urlprefix\url{https://doi.org/10.2991/978-94-91216-35-0}

\bibitem{BaloghRudin1992}
Z.~Balogh, M.~E. Rudin, Monotone normality, Topology and its Applications
  47~(2) (1992) 115--127.

\bibitem{Lutzer1971}
D.~J. Lutzer, On generalized ordered spaces, Panstwowe Wydawnictwo Naukowe,
  1971.

\bibitem{Lutzer1969}
D.~J. Lutzer, A metrization theorem for linearly orderable spaces, Proceedings
  of the American Mathematical Society 22~(2) (1969) 557--558.

\bibitem{AkinHrbacek2002}
E.~Akin, K.~Hrbacek, {Complete homogeneous LOTS}, Topology Proc 26 (2001)
  367--406.

\bibitem{Banakh2008}
T.~Banakh, A homogeneous first-countable zero-dimensional compactum failing to
  be a left-topological group, Matematychni Studii 29~(2) (2008).

\bibitem{Collins1996}
P.~Collins, Monotone normality, Topology and its Applications 74~(1-3) (1996)
  179--198.

\bibitem{Shkarin2004}
S.~Shkarin, Monotonically normal topological vector spaces are stratifiable,
  Topology and its Applications 136~(1-3) (2004) 129--134.

\end{thebibliography}\end{document}

